\section{Ramification After Blowup Is Tame}\label{section:ramification_after_blowup_is_tame}
Throughout this section let $C$ be a projective smooth geometrically connected curve over $k$.
Let $\mdls = \sum n_i P_i$ be a fixed modulus. Let $U=C \setminus \mdls$.
For any modulus $\ndls \subset \mdls$, we define $Z_\ndls$ as the closed subscheme of $C^{(\deg \ndls)} = \Sym_{k}^{\deg \ndls}(C)$ 
defined by $\ndls$ as a point of $C^{(\deg \ndls)}$
We then define $X_\ndls$ as the blowup of $C^{(\deg \ndls)}$ at $Z_\ndls$,
and we denote by $E_\ndls = Z_\ndls \times_{C^{(d)}} X_\ndls$ the exceptional divisor of this blowup, it is irreducible of codimension 1.
We denote by $\eta_\ndls$ the generic point of $E_\ndls$. 

Diagrammatically:
\begin{equation}\label{equation:blowup_diagram_ndls}
  \begin{tikzcd}
\overline{\{\eta_\ndls \}} =  E_\ndls \arrow[r, hook] \arrow[d] & X_\ndls \arrow[d, "\pi_\ndls"] \\
Z_\ndls \arrow[r, hook]           & C^{(\deg \ndls)}          
\end{tikzcd}   
\end{equation}

Our main goal in this section is to prove \Cref{theorem:SymmetricPowerOfSheavesIsTamelyRamifiedReduction}:
\begin{theorem}\label{theorem:torsors_after_blowup_is_tame}
    Let $G$ be a finite abelian group. And let $\cP \to U$ be a $G$-torsor in $U_{\text{\'Et}}$ with ramification bounded by $\mdls = \sum n_i P_i$.
    Let $\ndls = n_i P_i \subset \mdls$ be a sub modulus of $\mdls$. Then $\cP^{[\deg \ndls]}$ is tamely ramified at $\eta_\ndls$.
\end{theorem}

Our approach proceeds in three stages. 
First, we establish the result for the cyclic case 
$G = \mathbb{Z}/p\mathbb{Z}$ and for groups $G$ where $\gcd(|G|, p) = 1$.
Next, we extend the proof to cyclic $p$-groups $G = \mathbb{Z}/p^r\mathbb{Z}$. 
Finally, for a general finite abelian group $G$, we conclude by applying the structure theorem for finite abelian groups 
and \Cref{prop:quotient_torsors} to decompose the torsor into these constituent cases.

\subsection{\texorpdfstring{Ramification of $G$-Torsors}{Ramification of G-Torsors}}

Generally speaking, assume we are given a locally Noetherian scheme $X$, 
a dense open $U \subset X$, a finite \'etale morphsm $f: Y \to U$ and a prime divisor $Z \subset X$ such that $Z \cap U = \emptyset$
and the local ring $\Oo_{X, \xi}$ of $X$ at the generic point $\xi$ of $Z$ is a discrete valuation ring. 
Setting $K_\xi = \Frac(\Oo_{X, \xi})$ we obtain a cartesian square
\[ \xymatrix{ \mathop{\mathrm{Spec}}(K_\xi ) \ar[r] \ar[d] & U \ar[d] \\ \mathop{\mathrm{Spec}}(\mathcal{O}_{X, \xi }) \ar[r] & X } \]
of schemes. 
In particular, we see that $Y \times _ U \mathop{\mathrm{Spec}}(K_\xi )$ is the spectrum of a finite separable algebra $L_\xi /K_\xi $. 
\begin{definition}\label{definition:tamely_ramified_in_codim_1}
We say: 
\begin{enumerate}
    \item $Y$ is \textit{unramified} over $X$ at $(Z, \xi)$ resp. $Y$ is \textit{tamely ramified} over $X$ at $(Z, \xi)$
        if $L_\xi /K_\xi $ is unramified, resp. tamely ramified with respect to $\mathcal{O}_{X, \xi }$
        (\Cref{definition:tamely_ramified_dvrs})
    
    \item $Y$ is unramified over $X$ in codimension $1$, resp. $Y$ is tamely ramified over $X$ in codimension $1$ if $L_\xi /K_\xi $ is unramified, resp. tamely ramified 
    with respect to $\mathcal{O}_{X, \xi }$ for every $(Z, \xi )$ as above. 
\end{enumerate}

More precisely, we decompose $L_\xi $ into a product of finite separable field extensions of $K_\xi $ and we require each of these to be unramified, resp. tamely ramified with 
respect to $\mathcal{O}_{X, \xi }$.

\end{definition}

Now back to our orginial assumptions, $C$ is projective smooth geometrically connected curve
over a field $k$ and $U=C\setminus \mdls$. Every $(Z, \xi)$ is now a closed point $P$ and $\Oo_{C, P}$ is a discrete valuation ring. 
Assume $f: Y \to U$ is actually $G$-torsor $\cP \to U$ in $U_{\text{\'Et}}$ for $G$ finite abelian group.

By passing to the completion $\widehat{\Oo}_{C,P}$, we obtain the complete local field $\widehat{K}_P$.
Restricting the torsor $\cP$ to the punctured formal neighborhood $\Spec(\widehat{K}_P)$
 allows us to study the ramification in a strictly local context. 
 This restriction yields a $G$-torsor over a field, which necessarily decomposes into a disjoint union of spectra of finite separable field extensions
 $$ \cP|_{\Spec(\widehat{K}_P)} \cong \bigsqcup \Spec(F) $$
 such that:
 \begin{enumerate}
    \item These extensions $F/\widehat{K}_P$ are all isomorphic and Galois. 
    \item The Galois group $H = \Gal(F/\widehat{K}_P)$ identifies with the stabilizer of any connected component of the pullback $\cP|_{\Spec(\widehat{K}_P)}$, 
     which is a subgroup of $G$. 
    \item  The number of such components is then given by the index $[G:H]$.
 \end{enumerate}

 \begin{definition}\label{definition:local_ramification_boundness}
    We say that the extension $F/\widehat{K}_P$ \textit{has ramification bounded by $r$} if the ramification group $H^r$ is trivial. 
    We say \textit{$\cP$ has ramification bounded by $r$ at $P$} if the corresponding local field extensions $F/\widehat{K}_P$  
    satisfy this condition.
 \end{definition}


\begin{definition}
    Let $\mdls = \sum n_i P_i$ be a modulus on $C$. We say the $G$-torsor $\cP \to U$ has 
    \textbf{ramification bounded by $\mdls$} if, for each point $P_i$, the local restriction 
    $\cP|_{\Spec(\widehat{K}_{P_i})}$ has ramification bounded by the coefficient $n_i$.
\end{definition}

%%%%%%%%%%%%%%%%%%%%%%%%%%%%%%%%%%%%%%%
This local property remains invariant under base change to the separable closure.
Let $\bar{s} = \Spec(\bar{k}) \to \Spec(k)$ 
be a geometric point. 
The higher ramification groups for $\cP|_{\Spec(\widehat{K}_P)}$ are isomorphic to those of 
the geometric restriction $\cP_{\bar{k}}$ at any point $\bar{x}$ lying over $P$. 
Thus an equivalent defintion:
\begin{definition}
    Let $\mdls = \sum n_i P_i$ be a modulus on $C$. 
    The $G$-torsor $\cP \to U$ has \textbf{ramification bounded by $\mdls$} 
    if for every geometric point $\bar{x}$ of $\mdls$ (with image $\bar{s}$ in $\Spec(k)$), 
    the restriction of $\cP$ to$$ \Spec(\widehat{\mathcal{O}}_{C_{\bar{k}},\bar{x}}) \times_{C_{\bar{k}}} U_{\bar{k}} $$
    has ramification bounded by the multiplicity of $\mdls_{\bar{s}}$ at $\bar{x}$ in the sense of 
    \Cref{definition:local_ramification_boundness}.
\end{definition}
These definitions are equivalent because $\widehat{\mathcal{O}}_{C_{\bar{k}},\bar{x}} \cong \widehat{\mathcal{O}}_{C,P} \otimes_k \bar{k}$. Furthermore, the notions of tame ramification and unramifiedness derived here coincide with the codimension-$1$ criteria in \Cref{definition:tamely_ramified_in_codim_1}.

There is another equiviliant formulation through the language of characters.
The group $G$, being a finite abelian group over $k$, is \'etale. 
Consequently, the correspondence between $G$-torsors and the fundamental group allows us to describe $G$-torsors
 via characters. Specifically, there is an isomorphism: (\Cref{cor:abelian_equivilance_fundamental_torsors})
\[ \mathrm{Hom}_{\mathrm{cont.}}(\pi_1^{\text{\'et}}(U, \overline{x}), G) \xrightarrow{\cong} \mathrm{Tors}(U_{\text{\'et}}, G) \]

In our local setting, where we consider the restriction to the complete local field $\widehat{K}_P$, 
the fundamental group identifies with the absolute Galois group 
$G_{\widehat{K}_P} := \mathrm{Gal}({\widehat{K}_P}^{\mathrm{sep}}/{\widehat{K}_P})$. 
Thus, the restricted $G$-torsor $\mathcal{P}|_{\mathrm{Spec}(\widehat{K}_P)}$ corresponds to a unique continuous homomorphism:
\[ \rho : G_{\widehat{K}_P} \to G \]

Under this correspondence, the torsor $\mathcal{P}$ has ramification bounded by $r$ at $P$ 
if and only if the homomorphism $\rho$ trivializes the $r$-th ramification group in the upper numbering:
\[ \rho(G_{\widehat{K}_P}^r) = \{ 1 \} \]

\subsubsection*{Basic Properties of Ramification of \texorpdfstring{$G$-Torsors}{G-Torsors}}
We now prove some elementary lemmas regarding the ramification of $G$-torsors.
\begin{lemma}\label{lemma:bounding_ramification_of_contracted_product}
    Let $G$ be a finite abelian group and $X$ be a locally Noetherian scheme over a field $k$. 
    Let $U \subset X$ be a dense open subset and let $(Z, \xi)$ be a prime divisor in the complement $X \setminus U$.
    Assume $\cP_1$ and $\cP_2$ are two $G$-torsors on $U_{\text{ét}}$, Such that $\cP_1$ has ramification bounded by $r_1$ at $(Z, \xi)$, and $\cP_2$ has ramification bounded by $r_2$ at $(Z, \xi)$.
    Then their product $\cP_1 \otimes^G \cP_2$ has ramification bounded by $max(r_1, r_2)$ at $(Z, \xi)$.
\end{lemma}
\begin{proof}
    Let $A=\widehat{\Oo_{X, \xi}}$ and let $K=Frac(A)$.
    Let $\rho_1, \rho_2: G_K \to G$ be the associated continious homomorphisms corresponding to the $G$-torsors 
    $\cP_1|_{\Spec(K)}$, $\cP_2|_{\Spec(K)}$.
    We finish by noticing that the associated character of $(\cP_1 \otimes^G \cP_2)|_{\Spec(K)}$ is $\rho=\rho_1 + \rho_2$.    
\end{proof}

\begin{lemma}\label{lemma:descend_ramification_along_etale}

Let $f: X \to Y$ be a morphism of locally Noetherian schemes. Let $U_X \subset X$ and $U_Y \subset Y$ be dense open subschemes such that $f^{-1}(U_Y) \subset U_X$. 

Let $(Z_X, \xi_X)$ and $(Z_Y, \xi_Y)$ be prime divisors of $X$ and $Y$ such that $Z_X \cap U_X = \emptyset$ and $Z_Y \cap U_Y = \emptyset$. 
Suppose $f(\xi_X) = \xi_Y$ and that $f$ is étale at $\xi_X$.
Let $\mathcal{P}$ be a $G$-torsor over $U_Y$, and let $f^{-1} \mathcal{P}$ be its pullback to $f^{-1}(U_Y) \subset U_X$. 
Then, the ramification of $f^{-1} \mathcal{P}$ at $\xi_X$ is bounded by $r$ if and only if the ramification of $\mathcal{P}$ at $\xi_Y$ is bounded by $r$.
\end{lemma}
\begin{proof}
The boundedness of ramification is determined by the behavior of the torsor over the completion of the local rings at the generic points. 
Let $A = \mathcal{O}_{Y, \xi}$ and $B = \mathcal{O}_{X, \xi_X}$ be the discrete valuation rings at the generic points, with fraction fields 
$K$ and $L$ respectively. Since $f$ is étale at $\xi_X$, the map $A \to B$ is a flat, unramified (hence weakly unramified) local homomorphism. 
Consequently, the extension of completions $\widehat{L} / \widehat{K}$ is a finite unramified extension of complete discretely valued fields.
We finish by noticing that the upper numbering filtration on the absolute Galois group is compatible with unramified base change.\footnote{
 Specifically, let $G_K = \operatorname{Gal}(K^{sep}/K)$ and $G_L = \operatorname{Gal}(L^{sep}/L)$. For an unramified extension, the Herbrand function is the identity, which implies that for any $r \geq 0$:
\[ G_L^r = G_K^r \cap G_L \]}
\end{proof}


Let $C$ be a smooth, proper (projective) curve over a field $k$, and let $U \subset C$ be a dense open subset.
Let $G$ be a finite abelian group. and let $\pi_U: U' \to U$ be a $G$-torsor in $U_{\text{\'Et}}$. 
Let $C'$ be the normal proper model of $U'$ over $k$ (i.e. the unique smooth proper curve which has $U'$ as a dense open subset).
Let $\pi: C' \to C$ be the induced morphism. (which is finite generically \'etale over $U$). 
Let $U'^{(d)}$ be the standard $d$-th symmetric product of $U'$ as a scheme, and let $U'^{[d]}$ be the torsor-theoetric symmetric product,
by \Cref{cor:sym-power-quotient} we have a canonical morphism $q_U: U'^{(d)} \to U'^{[d]}$.
Note that $C'^{(d)}$ is smooth and finite over $C^{(d)}$, hence it is the normalization of $C^{(d)}$ inside $K(C'^{(d)})$.
$U'^{(d)}$ is a dense open subset of $C'^{(d)}$, hence $C'^{(d)}$ is a normal proper model of $U'^{(d)}$.
Let $C'^{[d]}$ be the normalization of $C^{(d)}$ in the function field $K(U'^{[d]})$,
then $q_U$ extends to a map $q: C'^{(d)} \to C'^{[d]}$, and we have the following diagrams:
\[
\begin{tikzcd}
	U' \arrow[r, hook] \arrow[d, "\pi_U"'] & C' \arrow[d, "\pi"] \\
	U \arrow[r, hook] & C
\end{tikzcd}
\qquad \qquad
\begin{tikzcd}
	{U'^{(d)}} \arrow[r, hook] \arrow[d, "q_U"'] & {C'^{(d)}} \arrow[d, "q"] \\
	{U'^{[d]}} \arrow[r, hook] \arrow[d] & {C'^{[d]}} \arrow[d] \\
	{U^{(d)}} \arrow[r, hook] & {C^{(d)}}
\end{tikzcd}
\]

\begin{lemma}\label{lemma:torsor_unramified_codim_one}
In the context of the preceding discussion, the morphism $U'^{(d)} \to U'^{[d]} \hookrightarrow C'^{[d]}$
is unramified in codimension $1$ (see \Cref{definition:tamely_ramified_in_codim_1}).
\end{lemma}
\begin{proof}

Let $C' \setminus U' = \{x_1, \ldots, x_m \}$. For each $x \in C' \setminus U'$, define the locus 
$$Z_x = \{D \in C'^{(d)} \mid x \in \mathrm{Supp}(D) \}.$$
Each $Z_x$ is a prime divisor in the symmetric product $C'^{(d)}$, and the boundary $C'^{(d)} \setminus U'^{(d)}$ is precisely the finite union $\bigcup_{i=1}^m Z_{x_i}$. 

Let $\eta$ be the generic point of $Z_x$. 
It suffices to show that the local ring extension $\mathcal{O}_{C'^{[d]}, q(\eta)} \hookrightarrow \mathcal{O}_{C'^{(d)}, \eta}$ is weakly unramified (\Cref{definition:weakly_unramified_and_tamely_ramified_DVR}).

Let $p: C'^d \to C'^{(d)}$ be the natural quotient. Consider the closed subvariety
$$H = \{x\} \times C' \times \cdots \times C' \subset C'^d.$$

The subvariety $H$ is irreducible, and its image under $p$ is exactly $Z_x$. 
Consequently, $p$ maps the generic point $\eta_H$ of $H$ to the generic point $\eta$ of $Z_x$.
\[
\begin{tikzcd}
\eta_H \arrow[d] \arrow[r] & H \arrow[r] \arrow[d] & C'^d \arrow[d, "p"] \\
\eta \arrow[r] \arrow[d] & Z_x \arrow[r] \arrow[d] & C^{(d)} \arrow[d, "q"] \\
q(\eta) \arrow[r] & q(Z_x) \arrow[r] & C'^{[d]}
\end{tikzcd}
\]
It suffices to show that the local ring extension $\mathcal{O}_{C'^{[d]}, q(\eta)} \hookrightarrow \mathcal{O}_{C'^d, \eta_H}$ is weakly unramified.
Let $R = \mathcal{O}_{C'^d, \eta_H}$. And let $A=\mathcal{O}_{C'^{[d]}, q(\eta)}$.

Recall that $C'^{[d]}$ is the quotient of $C'^d$ by the finite group $\Xi = \KG \rtimes S_d$ (see \Cref{cor:sym-power-quotient}),
and the morphism $q \circ p : C'^d \to C'^{[d]}$ is \textit{branched Galois Cover} (finite, surjective map of normal integral schemes $X \to Y$, with function field extension $K(X)/K(Y)$ Galois).

Hence the decomoposition group of $\eta_H$ in $\Xi$ and the inertia group of $\eta_H$ in $\Xi$ are defined:
$$D_{\eta_H} = \{ g \in \Xi \mid g(\eta_H) = \eta_H \}$$
$$I_{\eta_H} = \{ g \in D_{\eta_H} \mid g^*(u) \equiv u \pmod{\mathfrak{m}_{R}} \text{ for all } u \in R \}$$

%Then $A=R^{\Xi}$ where $\Xi = \KG \rtimes S_d$\footnote{Recall $\KG = \{ (g_1, \dots, g_d) \in G^d \mid g_1 g_2 \dots g_d = 1_G \}$} is the finite group acting on $C'^d$ as in \Cref{cor:sym-power-quotient}.
\textbf{Claim: } The inertia group $I_{\eta_H}$ is trivial.

Let $K=K(C')$ be the function field of $C'$, then 
$K(C'^d) = K^{\otimes_k d} \cong Frac(K_1 \otimes_k K_2 \otimes_k \cdots \otimes_k K_d)$
Where $K_i$ is the copy of $K$ corresponding to the $i$-th factor of $C'^d$.
The closed subscheme $H$ corresponds to the prime ideal $\pp_H \subset \Oo_{C'^d}$, and $R=\Oo_{C'^d, \pp_H}$. 
The residue field of $R$ is exactly $K(H) \cong Frac(k(x) \otimes_k K_2 \otimes_k \cdots \otimes_k K_d)$ where $k(x)$ is the residue field of $x$.

The action of $\gamma = (g_1, \ldots, g_d, \sigma) \in \Xi$, via the pullback $(\gamma^*)$ on $K(C'^d)$ does two things (in that order):
\begin{enumerate}
    \item It permutes the tensor factors according to $\sigma$. A function strictly in $K_i$ is mapped to a function in $K_{\sigma(i)}$.
    \item It applies the field automorphism $g_i^* \in \mathrm{Aut}(K/k)$ to the respective factors.
\end{enumerate}
Let $\gamma \in I_{\eta_H} \subset \Xi$, Then the induced automorphism $\bar{\gamma}^*$ on $K(H)$ must be the identity map.
This implies that $\sigma$ must fix $i$ for every $i \geq 2$ (The $K_i$ are linearly independent), and hence $\sigma=id$.
The trivial action of $\bar{\gamma}^*$ restricts to a trivial action of $g_i^*$ on $K_{i}$ for every $i \geq 2$, and hence 
$g_i = 1$ for every $i \geq 2$ (since the action of $G$ on $K(C')$ is faithful).
Finally, since $g_1 g_2 \cdots g_d = 1$, we also have $g_1 = 1$.
Since the inertia group $I_{\eta_H}$ is trivial, the local ring extension $A \hookrightarrow R$ is weakly unramified.
(See for example \stackstag{09E3}, \stackstag{0BTF})

\end{proof}


                % Let $\eta$ be a generic point of an irreducible component of the boundary divisor $C'^{(d)} \setminus U'^{(d)}$.
                %  Does the rational map $\Phi: C'^{(d)} \dashrightarrow C'^{[d]}$ extend to a morphism in an open neighborhood of $\eta$, and if so, is the corresponding local ring extension $\mathcal{O}_{C'^{[d]}, \Phi(\eta)} \hookrightarrow \mathcal{O}_{C'^{(d)}, \eta}$ unramified?

                % % \[
                % % \begin{tikzcd}
                % % U'^{(d)} \arrow[r, "q_U"] \arrow[d] & U'^{[d]} \arrow[r] \arrow[d] & C'^{(d)} \arrow[d] \\
                % % \mathrm{Sym}_k^d(U) \arrow[r] & \mathrm{Sym}_k^d(U) \arrow[r] & \mathrm{Sym}_k^d(C)
                % % \end{tikzcd}
                % % \]

                % The following lemma suggests that the map $q: U'^{(d)} \to U'^{[d]}$ from \Cref{cor:sym-power-quotient} is unramified in codimension 1 
                % according to \Cref{definition:tamely_ramified_in_codim_1}. 
                % While this is essentially true, the statement should not be taken literally in that specific context. 
                % Strictly adhering to \Cref{definition:tamely_ramified_in_codim_1} would require choosing a compactification of $U'^{[d]}$, 
                % which we prefer to avoid. 
                % Instead, we provide the following formulation, which suffices for our purposes and accurately characterizes the ramification behavior at codimension 1.

                % \begin{lemma}\label{lemma:symmetric_powers_of_torsors_and_ramification}
                    



                %      Then $\pi$ is unramified in codimension 1, i.e. for every prime divisor $(Z, \xi)$ in $\mathrm{Sym}_k^d(C)$ such that $Z \cap \mathrm{Sym}_k^d(U) = \emptyset$, the base change of $\pi$ to $\Spec(K_\xi)$ is unramified.
                %     Let $G$ be a finite abelian group.
                %     Assume we are given:
                %     \[
                %     \begin{tikzcd}
                %     U' \arrow[r] \arrow[d] & C' \arrow[d] \\
                %     U \arrow[r] & C
                %     \end{tikzcd}
                %     \]
                %     Where $C$ and $C'$ are smooth curves over a perfect field $k$, $U \subset C$ and $U' \subset C'$ are dense open subsets,
                %     and $U' \to U$ is a $G$-torsor
                %     Let 
                %     \[
                %         \begin{tikzcd}
                %        \mathrm{Sym}_k^d(U') \arrow[rd] \arrow[r, "q"] & U'^{[d]} \arrow[r] \arrow[d, "\pi"] & \mathrm{Sym}_k^d(C') \arrow[d] \\
                %        & \mathrm{Sym}_k^d(U) \arrow[r] & \mathrm{Sym}_k^d(C)
                %         \end{tikzcd}
                %     \]
                %     be the associated symmetric product diagram where $q: \mathrm{Sym}_k^d(U') \to U'^{[d]}$ is the quotient map defined in \Cref{cor:sym-power-quotient}
                %     and $\pi: U'^{[d]} \to \mathrm{Sym}_k^d(U)$ is a $G$-torsor. 
                %     Let $(Z, \xi)$ be a prime divisor in $\mathrm{Sym}_k^d(C)$ such that $Z \cap \mathrm{Sym}_k^d(U) = \emptyset$,
                %     let $\Oo_{C, \xi}$ be the corresponding local ring, and let $K_\xi = \Frac(\Oo_{C, \xi})$ be its fraction field. 
                %     Then 
                %     \begin{enumerate}
                %         \item $q$ and $\pi$ are \'etale.
                %         \item 
                %     \end{enumerate}
                %     Then the base change of $q$ to $\Spec(K_\xi)$ is unramified, i.e. the morphism
                %     \[ q_{K_\xi}: \mathrm{Sym}_k^d(U') \times_{\mathrm{Sym}_k^d(U)} \Spec(K_\xi) \to U'^{[d]} \times_{\mathrm{Sym}_k^d(U)} \Spec(K_\xi) \]
                %     is unramified.
                % \textcolor{red}{complete tomorrow}
                % \end{lemma}
                % Generally speaking, assume we are given a locally Noetherian scheme $X$, 
                % a dense open $U \subset X$, a finite \'etale morphsm $f: Y \to U$ and a prime divisor $Z \subset X$ such that $Z \cap U = \emptyset$
                % and the local ring $\Oo_{X, \xi}$ of $X$ at the generic point $\xi$ of $Z$ is a discrete valuation ring. 
                % Setting $K_\xi = \Frac(\Oo_{X, \xi})$ we obtain a cartesian square
                % \[ \xymatrix{ \mathop{\mathrm{Spec}}(K_\xi ) \ar[r] \ar[d] & U \ar[d] \\ \mathop{\mathrm{Spec}}(\mathcal{O}_{X, \xi }) \ar[r] & X } \]
                % of schemes. 
                % In particular, we see that $Y \times _ U \mathop{\mathrm{Spec}}(K_\xi )$ is the spectrum of a finite separable algebra $L_\xi /K_\xi $. 

                % We say: 
                % \begin{enumerate}
                %     \item $Y$ is \textit{unramified} over $X$ at $(Z, \xi)$ resp. $Y$ is \textit{tamely ramified} over $X$ at $(Z, \xi)$
                %         if $L_\xi /K_\xi $ is unramified, resp. tamely ramified with respect to $\mathcal{O}_{X, \xi }$
                %         (\Cref{definition:tamely_ramified_dvrs})
                    
                %     \item $Y$ is unramified over $X$ in codimension $1$, resp. $Y$ is tamely ramified over $X$ in codimension $1$ if $L_\xi /K_\xi $ is unramified, resp. tamely ramified 
                %     with respect to $\mathcal{O}_{X, \xi }$ for every $(Z, \xi )$ as above. 
                % \end{enumerate}


    %As in \Cref{cor:sym-power-quotient}
%$\mathrm{Sym}_S^d(\cP) \to \cP^{[d]}$


% \subsection{The Local Ring \texorpdfstring{$\Oo_{\eta_\ndls}$}{at the generic point of the blowup}}
% We remind the settings of this section: $C$ is a projective smooth geometrically connected curve over $k$.
% $\mdls = \sum n_i P_i$ is a fixed modulus. $U=C \setminus \mdls$.
% $\ndls \subset \mdls$ is a submodulus, $Z_\ndls$ is the closed subscheme of $C^{(\deg \ndls)} = \Sym_{k}^{\deg \ndls}(C)$ 
% defined by $\ndls$ as a point of $C^{(\deg \ndls)}$.
% $X_\ndls$ is the blowup of $C^{(\deg \ndls)}$ along $Z_\ndls$, $E_\ndls = Z_\ndls \times_{C^{(d)}} X_\ndls$ 
% is the exceptional divisor of this blowup, and $\eta_\ndls$ is the generic point of $E_\ndls$. 

% Diagrammatically:
% \begin{equation}\label{equation:blowup_diagram_ndls}
%   \begin{tikzcd}
% \overline{\{\eta_\ndls \}} =  E_\ndls \arrow[r, hook] \arrow[d] & X_\ndls \arrow[d, "\pi_\ndls"] \\
% Z_\ndls \arrow[r, hook]           & C^{(\deg \ndls)}          
% \end{tikzcd}   
% \end{equation}

% We now assume $P=P_i$, $n = n_i$ and $\ndls = n P$.

% Let $x \in \Oo_{C, P}$ be a local coordinate (uniformizer) at $P$.
% $C^{(\deg \ndls)}$ is a smooth variety of dimension $n$, 
% and $\ndls \in C^{(\deg \ndls)}$ is the point "serving as the origin" in this local context.

% Near $\ndls$, the coordinate ring of $C^{(\deg \ndls)}$ is generated by the elementary symmetric polynomials in 
% the local coordinates of the $n$ copies of $C$. If $x_1, \dots, x_n$ are the pullbacks of $x$ to $C^n$, 
% the local coordinates for $X$ at $\ndls$ are:
% $$e_1 = \sum_{i=1}^n x_i, \quad e_2 = \sum_{1 \le i < j \le n} x_i x_j, \dots, \quad e_n = x_1 x_2 \dots x_n$$


% Recall that $X_{\ndls} \subset C^{(\deg \ndls)} \times_k \mathbb{P}^{n-1}_k$ is defined locally by the equations $e_i u_j = e_j u_i$, where $[u_1 : \dots : u_n]$ are the homogeneous coordinates of $\mathbb{P}^{n-1}_k$.
% The exceptional divisor $E_{\ndls} \subset X_{\ndls}$ is the fiber over $\ndls$, $E_{\ndls} = \{ (\ndls, [u_1 : \dots : u_n]) \}$, which is isomorphic to $\mathbb{P}^{n-1}_k$. 
% Let $R = \mathcal{O}_{X_{\ndls}, \eta_{\ndls}}$. This ring $R$ is a discrete valuation ring (DVR) with fraction field $K = K(X) = k(C^n)^{S_n}$.
% Its residue field is  $\kappa(\eta_{\ndls}) = k(\frac{e_2}{e_1}, \dots, \frac{e_n}{e_1})$, 
% and the completion of $R$ with respect to its maximal ideal is:
% $$\hat{R} = \kappa(\eta_{\ndls})[[e_1]]$$
% In this local ring, the first elementary symmetric polynomial $e_1$ is a uniformizer. Any other coordinate $e_i$ is also a uniformizer as $e_i = (\frac{u_i}{u_1}) e_1$ and $\frac{u_i}{u_1}$ is a unit in $R$.

% For any symmetric polynomial $f \in k[x_1, \dots, x_n]^{S_n} = \Oo{C^{(n), \ndls}}$, we can write $f$ as a polynomial $G(e_1, \dots, e_n)$. If we expand $G$ into its homogeneous parts with respect to the $e_i$ coordinates, $G = G_d + G_{d+1} + \dots$, where $G_i$ is homogeneous of degree $i$ in $(e_1, \dots, e_n)$, the valuation $\nu_E$ associated with $E$ satisfies:
% $$\nu_E(f) = d$$
% This $d$ corresponds to the order of vanishing of the symmetric function $f$ at the point $\ndls$ in the symmetric product.


% \subsubsection*{The Affine Case}
% Let $X = \mathbb{A}_k^n$ be the affine $n$-space over a field $k$, and let $0 \in X$ be the origin. 
% Let $\tilde{X} = \text{Bl}_0(X) = X_{\ndls}$ be the blowup of $X$ at the origin. 
% Recall that $\tilde{X} \subset X \times_k \mathbb{P}^{n-1}_k$ is defined by the equations $x_i u_j = x_j u_i$, where $[u_1 : \dots : u_n]$ are the homogeneous coordinates of 
% $\mathbb{P}^{n-1}_k$.
% The exceptional divisor $E \subset \tilde{X}$ is the fiber over the origin, $E = \{ (0, [u_1 : \dots : u_n]) \}$, which is of codimension 1 in $\tilde{X}$. 
% Let $\eta \in E$ be the generic point of $E$, and let $R = \mathcal{O}_{\tilde{X}, \eta}$ be the associated local ring. 
% This ring $R$ is a discrete valuation ring (DVR) with fraction field $K = K(X) = k(x_1, \dots, x_n)$.

% On the affine chart $U_1$ where $u_1 \neq 0$, we have $x_i = \frac{u_i}{u_1} x_1$. 
% The coordinate ring is:$$\mathcal{O}_{\tilde{X}}(U_1) = k\left[x_1, \frac{u_2}{u_1}, \dots, \frac{u_n}{u_1}\right]$$
% In this chart, the generic point $\eta$ corresponds to the prime ideal $\mathfrak{p}_1 = (x_1)$. Thus, the local ring is $R = k[x_1, \frac{u_2}{u_1}, \dots, \frac{u_n}{u_1}]_{(x_1)}$. 
% The residue field is $\kappa(\eta) = k(\frac{u_2}{u_1}, \dots, \frac{u_n}{u_1})$, and the completion of $R$ with respect to its maximal ideal is:
% $$\hat{R} = \kappa(\eta)[[x_1]]$$
% In this local ring, $x_1$ is a uniformizer. Note that any $x_i$ (for $i > 1$) can also serve as a uniformizer, 
% as $x_i = (\frac{u_i}{u_1}) x_1$ and $\frac{u_i}{u_1}$ is a unit in $R$.

% For any monomial $M = x_1^{a_1} \dots x_n^{a_n} \in K$, we can write:
% $$M = x_1^{\sum a_i} \left( \frac{u_2}{u_1} \right)^{a_2} \dots \left( \frac{u_n}{u_1} \right)^{a_n}$$
% Since the term in parentheses is a unit in $R$, the valuation $\nu_E$ associated with $E$ satisfies:
% $$\nu_E(M) = \sum a_i = \deg M$$
% Consequently, for any polynomial $f = f_d + f_{d+1} + \dots + f_l$, where $f_i$ is the homogeneous part of degree $i$, we have $\nu_E(f) = d$ (the order of vanishing at the origin).

%\subsubsection*{The Symmetric Product Case}

% \textcolor{red}{in the above replace $n$ with $d$}
% \subsection{Proof of Theorem \ref{theorem:SymmetricPowerOfSheavesIsTamelyRamifiedReduction}}
% We assume $k$ is algebraically closed (We can base change by the unramified $\Spec{(k^{sep})} \to \Spec{(k)}$)
% We begin with:

% \begin{theorem}
%     Assume $G=\Z/p\Z$ and let $\cP \to U$ be a $G$-torsor which is wildly ramified at $P$ with ramification bounded by $d$.
%     Let $\ndls = dP$, the $G$-torsor $\cP^{[d]} \to U^{(d)}$ is unramified at $\eta_\ndls$- the generic point of 
%     $E_\mdls \subset X_{\mdls}$.
% \end{theorem}
% \begin{proof}
%     In order to be constructive, and present simpler proofs, we assume that 
%     $C=\bP_k^1$, $\mdls = d \cdot 0$. $U$

%     The computation that we will present is formal--local at $P$, , 
%     Once you replace $\mathbf{P}^1$ by an arbitrary smooth curve $C$, 
%     nothing essential changes at the generic point of the exceptional divisor, 
%     because (after choosing a uniformizer $x$ at $P$) the completed local ring $\widehat{\mathcal{O}}_{C,P}$ 
%     is still $k[[x]]$, and all the ``$d$--fold / symmetric / blowup'' constructions you use 
%     are compatible with \'etale (indeed formal) localization.

%     the general result will then be by 
% $\bG_m = U \subset U' = C \setminus \mdls$
%     Let $x$ be a local coordinate for $\Oo_{C, P}$, then 
%     $R=\widehat{\Oo_{C, P}} = k[[x]]$ and $K=\Frac{R}=k((x))$
%     Since $\Z/p\Z$ is simple, $\cP$ is either connected or a finite union of identity morphisms. The latter case being redundant. 
%     So, by \Cref{theorem:artin_schreier_ramification} $\cP|_{K} = \Spec(F)$, 
%     Where $F=\frac{K[X]}{(X^p - X - f(x))}$ and $f(x) = c x^{-m} + a_{-m+1}x^{-m+1} + \dots a_{-1}x^{-1} + a_0 \in k((x))$ with $m < d$.
%     Thus, The ring corresponding to $\widehat{O_{C^{(d)}, (P, \dots P)}}$ is 
%     $R' = k[[x_1, \dots, x_n]]$ with fraction field $K' = k((x_1, \dots, x_n))$.
%     And the pullback of field.    $p_1^{-1}(\cP) \times \dots \times p_d^{-1}(\cP)$ on $C^{d}$
%     Is 



%     In the previous section, we saw that when $x$ is local coordinate for $\Oo_{C, P}$ (a uniformizer) at $P$. 
%     Then $$\hat{R} = k(\frac{e_2}{e_1}, \dots, \frac{e_n}{e_1})[[e_1]]$$ where $R = \mathcal{O}_{X_{\ndls}, \eta_{\ndls}}$.
%     Thus if $K=Frac(R)$ then:
%     \begin{equation}\label{eq:complete_field_at_generic_point}
%           \hat{K}= k(\frac{e_1}{e_d}, \dots, \frac{e_{d-1}}{u_d})((e_d))
%      \end{equation}
    
    
% \end{proof}



\subsection{Proof of Theorem \ref{theorem:SymmetricPowerOfSheavesIsTamelyRamifiedReduction}}

We argue that it is enough to prove \Cref{theorem:torsors_after_blowup_is_tame} for $C=\bP^1_k$, $\ndls = d \cdot 0$, and $\cP \to \bG_m = \bP^1_k \setminus \{0, \infty\}$ a $G$-torsor given by an explicit global equation.
Informally, this is because \emph{the ramification of $\cP^{[d]}$ at $\eta_\ndls$ is a formal--local invariant of $\cP$ at $P$}:
it depends only on the completed local ring $\widehat{\Oo_{C,P}} \cong k[[x]]$ together with the restriction of $\cP$
to the punctured formal disc $\Spec(k((x)))$, since all of the constructions involved --- self products, quotients by finite
groups, and blowups --- commute with the flat base change to the completion.
We emphasize that this is not a general principle that can simply be quoted: it is a statement about the particular
constructions at hand, and it is proved by checking each construction separately. This is the content of
\Cref{lemma:formal_local_invariance} and \Cref{cor:reduction_to_P1} below.
The second ingredient of the reduction is \Cref{lemma:realization_over_Gm}: every local torsor class appearing in
\Cref{theorem:torsors_after_blowup_is_tame} is realized by an explicit torsor over $\bG_m \subset \bP^1_k$.

Throughout this subsection we assume that $k$ is algebraically closed; in particular every closed point of $C$ is $k$-rational, and $\ndls = dP$ with $P \in C(k)$ and $d = \deg \ndls$.

\begin{rem*}
    The assumption that $k$ is algebraically closed is harmless in our situation, for the following reasons.
    \begin{enumerate}
        \item \textbf{(Base change to $\bar{k}$.)} Assume $k$ is perfect, so that $\bar{k}=k^{sep}$.
        The formation of $C^{(d)}$ (a quotient by a finite group), of the blowup $X_\ndls$ along $Z_\ndls$, and of
        $\cP^{[d]}$ (\Cref{cor:sym-power-quotient}) commutes with the flat base change $\Spec(\bar{k}) \to \Spec(k)$.
        Since $E_\ndls \cong \bP^{\deg \ndls - 1}_k$ is geometrically irreducible, $(E_\ndls)_{\bar{k}}$ is irreducible
        and its generic point $\bar{\eta}$ lies over $\eta_\ndls$; moreover
        $\Oo_{X_\ndls, \eta_\ndls} \to \Oo_{(X_\ndls)_{\bar{k}}, \bar{\eta}}$ is a weakly unramified extension of discrete
        valuation rings with separable residue field extension. The proof of \Cref{lemma:descend_ramification_along_etale}
        uses only these two properties of the extension (there the extension was moreover finite étale), so it applies
        verbatim: the ramification of $(\cP_{\bar{k}})^{[d]} = (\cP^{[d]})_{\bar{k}}$ at $\bar{\eta}$ is bounded by $r$ if
        and only if the ramification of $\cP^{[d]}$ at $\eta_\ndls$ is. Hence unramifiedness and tameness at $\eta_\ndls$
        may be checked after base change to $\bar{k}$.
        \item \textbf{(Rationality of $P$ --- a caveat.)} If $P$ is a closed point with $\kappa(P) \neq k$, then after base
        change to $\bar{k}$ the point $Z_{\ndls}$ corresponds to the divisor $\sum_{\bar{P} \mapsto P} n \bar{P}$ on
        $C_{\bar{k}}$, whose support consists of $[\kappa(P):k]$ distinct points, whereas the computation below treats a
        modulus supported at a single rational point. We therefore prove \Cref{theorem:torsors_after_blowup_is_tame} under
        the standing assumption that $k$ is algebraically closed; this suffices for our purposes, since in the proof of
        \Cref{theorem:SymmetricPowerOfSheafIsTamelyRamified} the field is assumed algebraically closed before the present
        theorem is invoked.
    \end{enumerate}
\end{rem*}

We now formulate the formal--local invariance precisely. Fix $P \in C(k)$ lying in the support of $\mdls$, set $\ndls = dP$,
and fix a uniformizer $x$ of $\Oo_{C,P}$; it induces an isomorphism $\widehat{\Oo_{C,P}} \cong k[[x]]$.
The scheme $\Spec(k[[x]])$ has two points --- the closed point, mapping to $P$, and the generic point, mapping to the generic
point of $C$ --- so the punctured formal disc $\Spec(k((x)))$ maps to $U$, and we denote by
$$\cP_x := \cP \times_U \Spec(k((x)))$$
the corresponding restriction of $\cP$; it is a $G$-torsor over $\Spec(k((x)))$.
Let $e_1, \dots, e_d \in k[[x_1, \dots, x_d]]$ be the elementary symmetric polynomials in $x_1, \dots, x_d$, and set
$$ V = \Spec\left(k[[e_1, \dots, e_d]][e_d^{-1}]\right), \qquad W = \Spec\left(k[[x_1, \dots, x_d]][(x_1 \cdots x_d)^{-1}]\right).$$
The inclusion $k[[e_1, \dots, e_d]] \subset k[[x_1, \dots, x_d]]$ induces a finite flat morphism $\hat{r}: W \to V$
(note that inverting $e_d = x_1 \cdots x_d$ inverts each $x_i$), and for each $i$ we let
$$ q_i : W \to \Spec(k((x))), \qquad x \mapsto x_i.$$
The group $S_d$ acts on $W$ over $V$ by permuting $x_1, \dots, x_d$, and $q_{\sigma(i)} = q_i \circ \sigma$.

\begin{lemma}[Formal--local invariance]\label{lemma:formal_local_invariance}
    In the above notation:
    \begin{enumerate}
        \item The choice of $x$ induces an isomorphism $\widehat{\Oo_{C^{(d)}, Z_\ndls}} \cong k[[e_1, \dots, e_d]]$ carrying
        the maximal ideal to $(e_1, \dots, e_d)$, such that the induced morphism $c: \Spec(k[[e_1, \dots, e_d]]) \to C^{(d)}$
        is flat and satisfies $c^{-1}(U^{(d)}) = V$; we write $c_V : V \to U^{(d)}$ for the restriction.
        \item There is an induced isomorphism $\widehat{\Oo_{X_\ndls, \eta_\ndls}} \cong \kappa[[e_d]]$, where
        $\kappa = k(\frac{e_1}{e_d}, \dots, \frac{e_{d-1}}{e_d})$. In particular
        $\widehat{K}_{\eta_\ndls} := \Frac(\widehat{\Oo_{X_\ndls, \eta_\ndls}}) \cong \kappa((e_d))$, and the canonical
        morphism $\Spec(\widehat{K}_{\eta_\ndls}) \to U^{(d)}$ factors as
        $\Spec(\kappa((e_d))) \to V \xrightarrow{c_V} U^{(d)}$, where the first map does not depend on $(C, \cP)$.
        \item There is a canonical isomorphism of $G$-torsors over $V$
        $$ \cP^{[d]} \times_{U^{(d)}} V \;\cong\; \Xi \backslash \left( q_1^{-1}\cP_x \times_W \cdots \times_W q_d^{-1}\cP_x \right).$$
    \end{enumerate}
    Consequently, the $G$-torsor $\cP^{[d]}|_{\Spec(\widehat{K}_{\eta_\ndls})}$ --- and hence whether the ramification of
    $\cP^{[d]}$ at $\eta_\ndls$ is bounded by $r$ (resp.\ tame, resp.\ unramified) --- depends only on $d$ and on the
    isomorphism class of the $G$-torsor $\cP_x$ over $k((x))$, and not otherwise on the pair $(C, \cP)$.
\end{lemma}

\begin{proof}
    \textbf{(1)} Since $P$ is a $k$-rational point of the smooth curve $C$, the uniformizer $x$ induces
    $\widehat{\Oo_{C,P}} \cong k[[x]]$. As completion at rational points takes products of $k$-varieties to completed tensor
    products, we get $\widehat{\Oo_{C^d, (P, \dots, P)}} \cong k[[x]] \widehat{\otimes}_k \cdots \widehat{\otimes}_k k[[x]] = k[[x_1, \dots, x_d]]$.
    Write $A = \Oo_{C^{(d)}, Z_\ndls}$ and $B = \Oo_{C^d, (P, \dots, P)}$.
    Since $k$ is algebraically closed, a closed point $(Q_1, \dots, Q_d) \in C^d$ maps to $Z_\ndls = dP$ under the finite
    projection $r: C^d \to C^{(d)}$ if and only if $Q_1 = \dots = Q_d = P$; hence $(P, \dots, P)$ is the unique point of
    $C^d$ over $Z_\ndls$, $B = (r_* \Oo_{C^d})_{Z_\ndls}$ is a finite $A$-module, and $A = B^{S_d}$ (formation of invariants
    commutes with localization at invariant primes).
    Since $\mathfrak{m}_B$ is the unique maximal ideal of $B$ over $\mathfrak{m}_A$, we have
    $\mathfrak{m}_B^N \subset \mathfrak{m}_A B \subset \mathfrak{m}_B$ for some $N$, so the $\mathfrak{m}_A$-adic and
    $\mathfrak{m}_B$-adic topologies on $B$ agree, and therefore $B \otimes_A \hat{A} \cong \hat{B} \cong k[[x_1, \dots, x_d]]$.
    Now $A = B^{S_d}$ is the kernel of $B \to \prod_{\sigma \in S_d} B$, $b \mapsto (\sigma(b) - b)_\sigma$; as $\hat{A}$ is
    flat over $A$, tensoring with $\hat{A}$ preserves this kernel, so
    $$\hat{A} \cong (B \otimes_A \hat{A})^{S_d} \cong k[[x_1, \dots, x_d]]^{S_d}.$$
    Finally, $k[[x_1, \dots, x_d]]^{S_d} = k[[e_1, \dots, e_d]]$: the ring $k[x_1, \dots, x_d]$ is a finite free module over
    $k[e_1, \dots, e_d]$, the same topological argument identifies $k[[x_1, \dots, x_d]]$ with
    $k[x_1, \dots, x_d] \otimes_{k[e_1, \dots, e_d]} k[[e_1, \dots, e_d]]$, and taking $S_d$-invariants --- again a kernel,
    preserved by the flat base change $k[e_1, \dots, e_d] \to k[[e_1, \dots, e_d]]$ --- yields $k[[e_1, \dots, e_d]]$.
    Under these identifications the maximal ideal of $\hat{A}$ is $(e_1, \dots, e_d)$, and $c$ is flat as the composition of
    the flat morphisms $\Spec(\hat{A}) \to \Spec(A) \to C^{(d)}$.

    It remains to check that $c^{-1}(C^{(d)} \setminus U^{(d)}) = V(e_d)$ as sets. We have
    $C^{(d)} \setminus U^{(d)} = \bigcup_{Q \in \mathrm{Supp}(\mdls)} Z_Q$ with $Z_Q = \{D : Q \in \mathrm{Supp}(D)\}$, as in
    the proof of \Cref{lemma:torsor_unramified_codim_one}. For $Q \neq P$ the divisor $Z_Q$ does not contain the point
    $Z_\ndls = dP$, so its ideal is not contained in $\mathfrak{m}_A$, and $c^{-1}(Z_Q) = \emptyset$. For $Q = P$, pull back
    along the finite surjective morphism $\Spec(\hat{B}) \to \Spec(\hat{A})$: the preimage of $Z_P$ in $C^d$ is
    $\bigcup_i p_i^{-1}(P)$, whose trace on $\Spec(\hat{B}) = \Spec(k[[x_1, \dots, x_d]])$ is
    $V(x_1) \cup \dots \cup V(x_d) = V(x_1 \cdots x_d) = V(e_d)$; by surjectivity, $c^{-1}(Z_P) = V(e_d)$ as claimed. Hence
    $c^{-1}(U^{(d)}) = V$, and likewise $\Spec(\hat{B}) \times_{C^d} U^d = W$.

    \textbf{(2)} Blowing up commutes with flat base change, and the center $Z_\ndls$ pulls back under $c$ to the closed
    point $V(e_1, \dots, e_d)$; hence
    $$\hat{Y} := \Bl_{(e_1, \dots, e_d)}\left(\Spec k[[e_1, \dots, e_d]]\right) \cong X_\ndls \times_{C^{(d)}} \Spec(\hat{A}),$$
    and the exceptional divisor $\hat{E} \subset \hat{Y}$ is the base change of $E_\ndls$ along
    $\Spec(\hat{A}/\hat{\mathfrak{m}}) = \Spec(\kappa(Z_\ndls))$, i.e.\ the projection $\mathrm{pr}: \hat{Y} \to X_\ndls$
    restricts to an isomorphism $\hat{E} \cong E_\ndls \cong \bP_k^{d-1}$.
    Let $\hat{\eta}$ be the generic point of $\hat{E}$. Then $\mathrm{pr}(\hat\eta) = \eta_\ndls$, and the induced local
    homomorphism of discrete valuation rings $\Oo_{X_\ndls, \eta_\ndls} \to \Oo_{\hat{Y}, \hat{\eta}}$ carries a uniformizer
    to a uniformizer (the ideal of $E_\ndls$ pulls back to the ideal of $\hat{E}$) and induces an isomorphism of residue
    fields (both are the function field of $\bP_k^{d-1}$); hence it induces an isomorphism on completions.
    The completion of $\Oo_{\hat{Y}, \hat{\eta}}$ is computed exactly as for the blowup of affine space at the origin, which
    we carry out in detail below: on the chart $u_d \neq 0$ we have $e_i = \frac{u_i}{u_d} e_d$, the local ring at
    $\hat\eta$ is $\left(\hat{A}[\frac{u_1}{u_d}, \dots, \frac{u_{d-1}}{u_d}]\right)_{(e_d)}$, its residue field is
    $\kappa = k(\frac{u_1}{u_d}, \dots, \frac{u_{d-1}}{u_d}) = k(\frac{e_1}{e_d}, \dots, \frac{e_{d-1}}{e_d})$, and its
    completion is $\kappa[[e_d]]$.
    Finally, $e_d$ is invertible in $\widehat{K}_{\eta_\ndls} = \kappa((e_d))$, so the composite
    $\Spec(\widehat{K}_{\eta_\ndls}) \to \Spec(\Oo_{\hat{Y}, \hat\eta}) \to \hat{Y} \to \Spec(\hat{A})$ factors through $V$,
    giving the asserted factorization of $\Spec(\widehat{K}_{\eta_\ndls}) \to U^{(d)}$ through $c_V$.

    \textbf{(3)} Let $h : \cP^{\times_k d} \to U^{(d)}$ denote the canonical morphism; it is finite, being the composition
    of the finite étale morphism $\cP^{\times_k d} \to U^d$ with the finite morphism $U^d \to U^{(d)}$. By
    \Cref{cor:sym-power-quotient},
    $\cP^{[d]} = \Xi \backslash \cP^{\times_k d} = \underline{\Spec}_{U^{(d)}}\left((h_* \Oo_{\cP^{\times_k d}})^{\Xi}\right)$.
    Pushforward along the affine morphism $h$ commutes with the flat base change $c_V : V \to U^{(d)}$, and the formation of
    $\Xi$-invariants is a kernel, hence also commutes with flat base change; therefore
    $$ \cP^{[d]} \times_{U^{(d)}} V \cong \Xi \backslash \left(\cP^{\times_k d} \times_{U^{(d)}} V\right).$$
    By step (1), $U^d \times_{U^{(d)}} V \cong W$, so
    $\cP^{\times_k d} \times_{U^{(d)}} V \cong \cP^{\times_k d} \times_{U^d} W$.
    For each $i$, the composition $W \to \Spec(k[[x_1, \dots, x_d]]) \to \Spec(k[[x]]) \to C$ --- the middle map being the
    completion of the projection $p_i : C^d \to C$, i.e.\ $x \mapsto x_i$ --- factors through $\Spec(k((x))) \to U$, since
    $x_i$ is invertible on $W$; this composite is precisely $q_i$, and it commutes with $p_i : U^d \to U$. Since
    $\cP^{\times_k d} = p_1^{-1}\cP \times_{U^d} \cdots \times_{U^d} p_d^{-1}\cP$, base changing to $W$ gives an isomorphism
    $$ \cP^{\times_k d} \times_{U^d} W \cong q_1^{-1}\cP_x \times_W \cdots \times_W q_d^{-1}\cP_x, $$
    which is equivariant for the actions of $G^d$ (in particular of $\KG$) and of $S_d$ --- the latter acting on $W$ by
    permuting the $x_i$, compatibly with $q_{\sigma(i)} = q_i \circ \sigma$. Passing to $\Xi$-quotients yields the displayed
    isomorphism.

    For the final statement: the schemes $V$ and $W$, the maps $q_i$ and $\hat{r}$, the $\Xi$-action, and the point
    $\Spec(\kappa((e_d))) \to V$ of (2) are all constructed from $d$ alone; by (3), the torsor
    $\cP^{[d]}|_{\Spec(\widehat{K}_{\eta_\ndls})}$ is the pullback along $\Spec(\kappa((e_d))) \to V$ of a torsor constructed
    from $\cP_x$ and these data. Since the ramification at $\eta_\ndls$ is by definition read off from
    $\cP^{[d]}|_{\Spec(\widehat{K}_{\eta_\ndls})}$ together with the discrete valuation ring $\kappa[[e_d]]$
    (\Cref{definition:tamely_ramified_in_codim_1}, \Cref{definition:local_ramification_boundness}), it depends only on
    $(\cP_x, d)$.
\end{proof}

\begin{cor}[Reduction to $\bP^1$]\label{cor:reduction_to_P1}
    Let $C, C'$ be smooth projective geometrically connected curves over the algebraically closed field $k$, let
    $\cP \to U = C \setminus \mdls$ and $\cP' \to U' = C' \setminus \mdls'$ be $G$-torsors, let
    $P \in \mathrm{Supp}(\mdls)(k)$ and $P' \in \mathrm{Supp}(\mdls')(k)$, and let $x, x'$ be uniformizers at $P, P'$
    respectively. Assume that $\cP_x \cong \cP'_{x'}$ as $G$-torsors, compatibly with the isomorphism
    $k((x)) \cong k((x'))$, $x \mapsto x'$.
    Then for every $d \geq 1$ and every $r$, the ramification of $\cP^{[d]}$ at $\eta_{dP}$ is bounded by $r$ if and only if
    the ramification of $\cP'^{[d]}$ at $\eta_{dP'}$ is; in particular, the one is unramified (resp.\ tamely ramified) at
    $\eta_{dP}$ if and only if the other is at $\eta_{dP'}$.
\end{cor}
\begin{proof}
    Immediate from \Cref{lemma:formal_local_invariance}: both $\cP^{[d]}|_{\Spec(\widehat{K}_{\eta_{dP}})}$ and
    $\cP'^{[d]}|_{\Spec(\widehat{K}_{\eta_{dP'}})}$ are identified, compatibly with the discrete valuation ring
    $\kappa[[e_d]]$, with the pullback along $\Spec(\kappa((e_d))) \to V$ of
    $\Xi \backslash (q_1^{-1}\cP_x \times_W \cdots \times_W q_d^{-1}\cP_x)$.
\end{proof}

\begin{lemma}[Realization over $\bG_m$]\label{lemma:realization_over_Gm}
    Let $k$ be algebraically closed of characteristic $p > 0$, and let $\cQ$ be a $G$-torsor over $\Spec(k((x)))$.
    \begin{enumerate}
        \item If $\cQ$ is tamely ramified, then there exist an integer $e$ prime to $p$, an injective homomorphism
        $\iota : \Z/e\Z \hookrightarrow G$, and an integer $a$ prime to $e$, such that
        $$\cP' := \iota_*\left( \Spec\left(k[x, x^{-1}][T]/(T^e - x^a)\right) \right)$$
        is a $G$-torsor over $\bG_m$, tamely ramified at $0$ and $\infty$, with $\cP'_x \cong \cQ$.
        Here $T^e = x^a$ is regarded as a $\Z/e\Z$-torsor via a fixed primitive $e$-th root of unity $\zeta_e \in k$, and
        $\iota_* \cQ' := (\cQ' \times G)/(\Z/e\Z)$ denotes the induced $G$-torsor.
        \item If $G = \Z/p\Z$ and $\cQ$ has ramification bounded by $d$ (\Cref{definition:local_ramification_boundness}),
        then there exists $f \in x^{-1}k[x^{-1}]$, either $f = 0$ or
        $f = cx^{-m} + a_{-m+1}x^{-m+1} + \dots + a_{-1}x^{-1}$ with $c \neq 0$, $p \nmid m$ and $m < d$, such that
        $$\cP' := \Spec\left(k[x, x^{-1}][X]/(X^p - X - f)\right)$$
        is a $\Z/p\Z$-torsor over $\bG_m$ with $\cP'_x \cong \cQ$. Moreover $\cP'$ extends to a torsor over
        $\bP^1_k \setminus \{0\}$ --- in particular it is unramified at $\infty$ --- and it has ramification bounded by $d$
        at $0$.
    \end{enumerate}
    In both cases $\cP'_x$ denotes the restriction of $\cP'$ to $\Spec(k((x)))$, as before.
\end{lemma}
\begin{proof}
    \textbf{(1)} Since $k$ is algebraically closed, the residue field of $k((x))$ is algebraically closed, so $k((x))$
    admits no non-trivial unramified extensions, and its tamely ramified extensions are exactly the Kummer extensions
    $k((x^{1/e}))$ with $p \nmid e$ (\Cref{theorem:kummer_ramification}). Hence the tame quotient $G^t$ of
    $\Gal(k((x))^{sep}/k((x)))$ is procyclic, $G^t \cong \varprojlim_{p \nmid e} \mu_e(k) \cong \prod_{\ell \neq p} \Z_\ell$;
    fix a topological generator $\gamma$.
    Let $\rho: \Gal(k((x))^{sep}/k((x))) \to G$ be the continuous homomorphism classifying $\cQ$. Tameness means that
    $\rho$ kills the wild inertia subgroup, so $\rho$ factors through $G^t$ and is determined by $g := \rho(\gamma)$, whose
    order $e$ is prime to $p$. Let $\iota: \Z/e\Z \xrightarrow{\sim} \langle g \rangle \subset G$ be given by
    $1 \mapsto g$, so that $\rho = \iota \circ \chi$ for a surjective character $\chi : G^t \to \Z/e\Z$.
    On the other hand, for $a \in \Z$ the equation $T^e = x^a$ defines a $\mu_e \cong \Z/e\Z$-torsor over
    $\Spec(k((x)))$, and by Kummer theory the classes of $x^a$, $a \in \Z/e\Z$, exhaust
    $\mathrm{Hom}_{\mathrm{cont.}}(G^t, \Z/e\Z) \cong \Z/e\Z$, the class of $x^a$ being surjective precisely when
    $\gcd(a, e) = 1$. Choose $a$ with the class of $x^a$ equal to $\chi$; then $\cQ \cong \iota_*(T^e = x^a)$ over
    $k((x))$. The same equation over $k[x, x^{-1}]$ defines a $\Z/e\Z$-torsor over $\bG_m$ (étale since $e$ is invertible
    in $k$), so $\cP' := \iota_*\left(\Spec(k[x,x^{-1}][T]/(T^e - x^a))\right)$ is a $G$-torsor over $\bG_m$ with
    $\cP'_x \cong \cQ$; it is tamely ramified at $0$ and at $\infty$ since its local characters there factor through the
    tame quotients.

    \textbf{(2)} If $\cQ$ is trivial, take $f = 0$. Otherwise, since $\Z/p\Z$ is simple and $k((x))$ admits no non-trivial
    unramified extensions, $\cQ = \Spec(F)$ for a cyclic extension $F/k((x))$ of degree $p$ together with an isomorphism
    $\Gal(F/k((x))) \cong \Z/p\Z$. By Artin--Schreier theory such data correspond to non-zero classes in
    $k((x))/\wp(k((x)))$, where $\wp(u) = u^p - u$, via $g \mapsto k((x))[X]/(X^p - X - g)$. We claim that every class
    admits a representative $f$ as in the statement:
    \begin{itemize}
        \item $\wp$ is surjective on $k[[x]]$: given $h = \sum_{i \geq 0} b_i x^i$, solve $\wp(u) = h$ with
        $u = \sum_{i \geq 0} u_i x^i$ by choosing $u_0$ with $u_0^p - u_0 = b_0$ (possible as $k$ is algebraically closed)
        and setting recursively $u_i = u_{i/p}^p - b_i$, with the convention $u_{i/p} = 0$ when $p \nmid i$. Hence, we may
        choose a representative in $x^{-1}k[x^{-1}]$.
        \item If the representative has pole order $m$ divisible by $p$, with leading term $cx^{-m}$, then subtracting
        $\wp(c^{1/p}x^{-m/p}) = cx^{-m} - c^{1/p}x^{-m/p}$ strictly decreases the pole order. Iterating, we reach a
        representative $f$ which is either $0$ (impossible, as the class is non-zero) or has pole order $m$ prime to $p$.
    \end{itemize}
    By \Cref{theorem:artin_schreier_ramification}, the unique ramification jump of $F = k((x))[X]/(X^p - X - f)$ is at $m$,
    so the ramification of $\cQ$ is bounded by $d$ if and only if $m < d$.
    Finally, $f \in k[x, x^{-1}]$, so the same equation defines a finite étale $\Z/p\Z$-cover $\cP' \to \bG_m$ (a torsor via
    $X \mapsto X + 1$) with $\cP'_x \cong \cQ$; and writing $y = x^{-1}$, we have $f \in y\,k[y]$, so $X^p - X - f$ defines
    a finite étale cover of $\mathbb{A}^1_k = \Spec(k[y]) = \bP^1_k \setminus \{0\}$ extending $\cP'$, whence $\cP'$ is
    unramified at $\infty$.
\end{proof}

Combining \Cref{cor:reduction_to_P1} with \Cref{lemma:realization_over_Gm}, in order to prove
\Cref{theorem:torsors_after_blowup_is_tame} for a torsor which is tamely ramified at $P$, or for a $\Z/p\Z$-torsor with
ramification bounded by $d$ at $P$, we may --- and do --- assume from now on that
$$C = \bP^1_k, \qquad \mdls = d \cdot 0 + \infty, \qquad \ndls = d \cdot 0, \qquad U = \bG_m,$$
and that $\cP \to \bG_m$ is one of the explicit torsors of \Cref{lemma:realization_over_Gm}.

We start by analyzing the blowup $\tilde{X} = \text{Bl}_0(X)$ of $X = \mathbb{A}_k^n$
where $0 \in X$ be the origin. 
Recall that $\tilde{X} \subset X \times_k \mathbb{P}^{n-1}_k$ is defined by the equations $x_i u_j = x_j u_i$, where $[u_1 : \dots : u_n]$ are the homogeneous coordinates of 
$\mathbb{P}^{n-1}_k$.
The exceptional divisor $E \subset \tilde{X}$ is the fiber over the origin, $E = \{ (0, [u_1 : \dots : u_n]) \}$, which is of codimension 1 in $\tilde{X}$. 
Let $\eta \in E$ be the generic point of $E$, and let $R = \mathcal{O}_{\tilde{X}, \eta}$ be the associated local ring. 
This ring $R$ is a discrete valuation ring (DVR) with fraction field $K = K(X) = k(x_1, \dots, x_n)$.

On the affine chart $U_1$ where $u_1 \neq 0$, we have $x_i = \frac{u_i}{u_1} x_1$. 
The coordinate ring is:$$\mathcal{O}_{\tilde{X}}(U_1) = k\left[x_1, \frac{u_2}{u_1}, \dots, \frac{u_n}{u_1}\right]$$
In this chart, the generic point $\eta$ corresponds to the prime ideal $\mathfrak{p}_1 = (x_1)$. Thus, the local ring is $R = k[x_1, \frac{u_2}{u_1}, \dots, \frac{u_n}{u_1}]_{(x_1)}$. 
The residue field is $\kappa(\eta) = k(\frac{u_2}{u_1}, \dots, \frac{u_n}{u_1})$, and the completion of $R$ with respect to its maximal ideal is:
$$\hat{R} = \kappa(\eta)[[x_1]]$$
In this local ring, $x_1$ is a uniformizer. Note that any $x_i$ (for $i > 1$) can also serve as a uniformizer, 
as $x_i = (\frac{u_i}{u_1}) x_1$ and $\frac{u_i}{u_1}$ is a unit in $R$.

For any monomial $M = x_1^{a_1} \dots x_n^{a_n} \in K$, we can write:
$$M = x_1^{\sum a_i} \left( \frac{u_2}{u_1} \right)^{a_2} \dots \left( \frac{u_n}{u_1} \right)^{a_n}$$
Since the term in parentheses is a unit in $R$, the valuation $\nu_E$ associated with $E$ satisfies:
$$\nu_E(M) = \sum a_i = \deg M$$
Consequently, for any polynomial $f = f_d + f_{d+1} + \dots + f_l$, where $f_i$ is the homogeneous part of degree $i$, we have $\nu_E(f) = d$ (the order of vanishing at the origin).



Our first result is:
\begin{theorem}\label{theorem:ramification_after_blowup_P1}
     Let $\cP \to \bG_m \subset \bP_k^1$ be a $G$ torsor which is either
     \begin{enumerate}
          \item tamely ramified at $0$ 
          \item totally ramified at $0$ with $G=\Z/p\Z$ and ramificaiton bounded by $d$.
     \end{enumerate}
     Then the ramification of the $G$-torsor $\cP^{[d]} \to \bG_m^{(d)}$ at $\eta_\mdls$ the generic point of $E_\mdls \subset X_{\mdls}$ is  
     \begin{enumerate}
          \item tamely ramified if $\cP$ was tamely ramified
          \item unramified if $\cP$ was totally ramified with $G=\Z/p\Z$ and ramification bounded by $d$
     \end{enumerate}
\end{theorem}
\begin{proof}
    The local ring at the generic point of the excptional divisor of the blowup of the affine space at 0 point is
     $R = k[x_d, \frac{u_1}{u_d}, \dots, \frac{u_{d-1}}{u_d}]_{(x_d)}$. 
     Where for every  $i < d$, we have $x_i = \frac{u_i}{u_d} x_d$ are all uniformizers.
     The residue field is $\kappa(\eta) = k(\frac{u_1}{u_d}, \dots, \frac{u_{d-1}}{u_d})$
     and the completion of $R$ with respect to its maximal ideal is:
     $$\hat{R} = \kappa(\eta)[[x_d]]$$
     In our situation, when we take symmetric product of the affine space, the situation is similiar with different coordinates
     if we let $e_1, \dots, e_d$ be the symmetric polynomials in $x_1, \dots x_d$ then:
     The local ring is 
     $R = k[e_d, \frac{u_2}{u_d}, \dots, \frac{u_{d-1}}{u_d}]_{(e_d)}$. 
     $e_i = \frac{u_i}{u_d} e_d$ are all uniformizers.
     The residue field being:
     $\kappa(\eta) = k(\frac{u_1}{u_d}, \dots, \frac{u_{d-1}}{u_d})$
     and the completion: $\hat{R} = \kappa(\eta)[[s_d]]$
     Note that from  $e_i = \frac{u_i}{u_d} e_d$
     We get  $\frac{e_i}{e_d} = \frac{u_i}{u_d}$ in the fracion field. hence

     \begin{equation}\label{eq:complete_field_at_generic_point}
          \hat{K}= k(\frac{e_1}{e_d}, \dots, \frac{e_{d-1}}{u_d})((e_d))
     \end{equation}

     We compute directly the extesntion of complete valued fields over the complete valued field at the generic point. 
     Note that by \Cref{cor:reduction_to_P1} and \Cref{lemma:realization_over_Gm}
     we can assume $\cP = \Spec k[x, x^{-1}][X]/(X^n - a)$ for $a \in k[x, x^{-1}]$ or $\cP = \Spec k[x, x^{-1}][X]/(X^p + X - f(x, x^{-1}))$
     where $f(x,x^{-1}) = c x^{-m} + a_{-m + 1} x^{-m + 1} + \cdots a_{-1}x^{-1} + a_0 = c x^{-m} + f_{-m+1}(x^{-1})$ where $m < d$.
     
     We deal with each case separately. 
     
     \textbf{Artin-Schreier Extensions:}

     Set $R=k[x, x^{-1}]$, and $S = \Spec k[x, x^{-1}][X]/(X^p + X - f(x^{-1}))$
     we have 
     $$
          \begin{aligned}
          R^{\otimes_k d} &= k[x_1, x_1^{-1}, \dots, x_d, x_d^{-1}] \\
          S^{\otimes_k d} &= k[x_1, x_1^{-1}, \dots, x_d, x_d^{-1}][X_1, \dots, X_d] / (X_1^p - X_1 - f(x_1), \dots, X_d^p - X_d - f(x_d)) \\
          &= R^{\otimes_k d}[X_1, \dots, X_d] / (X_1^p - X_1 - f(x_1), \dots, X_d^p - X_d - f(x_d))
          \end{aligned}
     $$

     Next, we want to understand the ring corresponding to $p_1^{-1}(\cP) \otimes \dots \otimes p_d^{-1}(\cP)$ on $C^{d}$ - 
     the $d$'th-product of the $G=\Z/p\Z$-torsors  $p_1^{-1}(\cP), \dots, p_d^{-1}(\cP)$ on $U^{d}$. 
     It corresponds to qoutient:
     $\left(p_1^{-1}(\cP) \times \dots \times p_d^{-1}(\cP)\right) / G^{d-1}$ where the action of $G^{d-1}$ on the product is:
     $$(g_1, \dots, g_{d-1}) \cdot (p_1, p_2, \dots, p_{d-1}, p_{d}) = (g_1(p_1), g_1^{-1}g_2(p_2), \dots, 
     g_{d-2}^{-1}g_{d-1}(p_{d-1}), g_{d-1}^{-1}(p_d))$$
     
     The affine ring corrsponding to the torsor product is $\left(S^{\otimes_k d}\right)^{G^{d-1}}$
    
     Recall that the action of $g \in G=\Z/p\Z$ on $X$ is $g(X) = X + g$ ($g$ correspond to a number $0 \leq g \leq p-1)$.
     So, the action of $(g_1, \dots, g_{d-1})$ on the generators $(X_1, X_2, \dots, X_{d-1}, X_{d})$
     Is $X_1 \mapsto X_1 + g_1$, $X_i \mapsto X_i - g_{i-1} + g_i$ for $1<i<d$ and $X_d \mapsto X_d - g_{d-1}$.
     So we see that $Y=X_1 + \dots + X_d$ is invariant.
     Moreover $Y^P - Y - \sum_{i=1}^d f(x_i)=0$ is irreducible degree $p$ equation for $Y$,
     Since we are quotienting a rank $p^{d}$ extesntion by a group of order $p^{d-1}$ the resulting invaraint subring must
     have rank $p$ over $R^{\otimes_k d}$, So we conclude:
     \[
          \left(S^{\otimes_k d}\right)^{G^{d-1}} \cong R^{\otimes_k d}[Y] / (Y^p - Y - \sum_{i=1}^d f(x_i))
     \]  
     
     The group $S_d$ acts on 
     $ R^{\otimes_k d} = k[x_1^{\pm 1}, x_2^{\pm 1}, \dots, x_d^{\pm 1}]$
      by permuting the variables $\{x_i\}_{i=1}^d$.
     And since $Y=\sum_i^d X_i$ it leaves $Y$ invaraint. 
     The invaraint subring $(R^{\otimes_k d})^{S_d}$ is simply $k[e_1, e_2, \dots, e_d, e_d^{-1}]$
     where $\{ e_i \}$  are the symmetric polynomials in $x_1, \dots, x_d$:
     $$e_k = \sum_{1 \le j_1 < j_2 < \dots < j_k \le d} x_{j_1} x_{j_2} \dots x_{j_k}$$
     
     To find $\left(R^{\otimes_k d}[Y] / (Y^p - Y - \sum_{i=1}^d f(x_i))\right)^{S_d}$ its enough to express
     $\sum_{i=1}^d f(x_i)$ in $e_1, \dots, e_d$, this can be done with the newton polynomials, moreover, we claim the following:

     \begin{lemma}
          Let $f(x) = c x^{-m} + a_{-m + 1} x^{-m + 1} + \cdots a_{-1}x^{-1} + a_0$, 
          define  $\alpha(x_1, ..., x_d) = \sum_1^d f(x_i)$, and deonte by $e_1, ..., e_d$ the elementary symmetric polynomials 
          in $x_1, ..., x_d$. If $m < d$ then $\alpha(x_1, ..., x_d) \in k(e_1/e_d, e_2/e_d, ..., e_{d-1}/e_d)$
     \end{lemma}
     \begin{proof}
          Changing variables $y_i=x_i^{-1}$ for each $i \in \{1, \dots, d\}$
          We get 
          $$\alpha = \sum_{i=1}^d f(x_i) = \sum_{i=1}^d \left( c y_i^m + a_{-m+1} y_i^{m-1} + \dots + a_{-1} y_i + a_0 \right)$$
          Rearranging the sums, we get
          $$\alpha = c \sum_{i=1}^d y_i^m + a_{-m+1} \sum_{i=1}^d y_i^{m-1} + \dots + a_{-1} \sum_{i=1}^d y_i + d a_0$$
          Let $p_k(y_1, \dots, y_d) = \sum_{i=1}^d y_i^k$ be the $k$-th power sum symmetric polynomial. 
          The expression for $\alpha$ is a linear combination of these power sums:
          $$\alpha = c p_m(y) + a_{-m+1} p_{m-1}(y) + \dots + a_{-1} p_1(y) + d a_0$$
         
          According to the \textit{Fundamental Theorem of Symmetric Polynomials}, any symmetric polynomial in $y_1, \dots, y_d$ 
          can be expressed as a polynomial in the elementary symmetric polynomials $e_k(y_1, \dots, y_d)$. 
          Since $m < d$, $\alpha$ is a polynomial in $e_1(y), e_2(y), \dots, e_m(y)$. ($y=(y_1, \dots, y_d)$)
          The elementary symmetric polynomials in $y_i = 1/x_i$ are related to the elementary symmetric polynomials in $x_i$ as follows: 
          $$e_k(y_1, \dots, y_d) = \sum_{1 \le i_1 < \dots < i_k \le d} \frac{1}{x_{i_1} \dots x_{i_k}} = \frac{\sum_{1 \le j_1 < \dots < j_{d-k} \le d} x_{j_1} \dots x_{j_{d-k}}}{x_1 x_2 \dots x_d}$$
          Thus,
          $$e_k(y_1, \dots, y_d) = \frac{e_{d-k}(x_1, \dots, x_d)}{e_d(x_1, \dots, x_d)}$$
          which concludes the proof.
     \end{proof}

Finally, restricting $\cP^{[d]}$ to $\spec {\hat{K}}$ we get by \Cref{eq:complete_field_at_generic_point}
and \Cref{theorem:artin_schreier_ramification} the result. (that $\cP$ is unramified at the generic point of the exceptional divisor of the blowup).

\textbf{Kummer Extensions:}
Few things are different in that case,

     Set $R=k[x, x^{-1}]$, and $S=R[X]/(X^n - f)$ where $f= f(x, 1/x) \in R$
     In this case we have $char k = p$ and $gcd(p, n)=1$.

     We have 
     $$
          \begin{aligned}
          R^{\otimes_k d} &= k[x_1, x_1^{-1}, \dots, x_d, x_d^{-1}] \\
          S^{\otimes_k d} &= k[x_1, x_1^{-1}, \dots, x_d, x_d^{-1}][X_1, \dots, X_d] / (X_1^n - f_1, \dots, X_d^n - f_d) \\
          &= R^{\otimes_k d}[X_1, \dots, X_d] / (X_1^n - f_1, \dots, X_d^n - f_d)
          \end{aligned}
     $$
     Where $f_i = f(x_i, x_i^{-1})$

     Next, we want to figure out 
      $\left(S^{\otimes_k d}\right)^{G^{d-1}}$
    

     Recall that the action of $g \in G=\Z/n\Z$ on $X$ is $g(X) = \zeta^{g} X$ ($g$ correspond to a number $0 \leq g \leq n-1)$.
     So, the action of $(g_1, \dots, g_{d-1})$ on the generators $(X_1, X_2, \dots, X_{d-1}, X_{d})$
     Is $X_1 \mapsto \zeta^{g_1} X_1$, $X_i \mapsto \zeta^{g_i - g_{i-1}}X_i$ for $1<i<d$ and $X_d \mapsto \zeta^{-g_{d-1}} X_d$.

     So we see that $Y=X_1 X_2 \dots X_d$ is invariant.
     And $Y^n - \prod_{i=1}^d f_i$ is irreducible degree $n$ equation for $Y$,
     So, like before, we conclude:
     \[
          \left(S^{\otimes_k d}\right)^{G^{d-1}} \cong R^{\otimes_k d}[Y] / (Y^n - \prod_{i=1}^d f_i)
     \]  
     
     The group $S_d$ acts on $R^{\otimes_k d}[Y] / (Y^n - \prod_{i=1}^d f_i)$ by permuting the indices.
     on the variables $x_i$, 
     On $Y=\prod_1^d X_i$ it is invaraint. 
     The invaraint subring $(R^{\otimes_k d})^{S_d}$ is simply $k[e_1, e_2, \dots, e_d, e_d^{-1}]$ like before.
     

     The polynomial $F=\prod_{i=1}^d f_i$ is symmetric in $\{x_i\}_1^d$ so it can be expressed as a polynomial
     $\tilde{F}(e_1, \dots, e_d)$ in the elementary symmetric variables. 
     Hence the qoutient ring is:
     $$\left( \frac{k[x_1^{\pm 1}, \dots, t_d^{\pm 1}][Y]}{(Y^n - \prod_{i=1}^d f_i)} \right)^{S_d} \cong 
     \frac{k[e_1, \dots, e_d, e_d^{-1}][Y]}{(Y^n - \tilde{F}(e_1, \dots, e_d))}$$
     So we see again, that  restricting $\cP^{[d]}$ to $\spec {\hat{K}}$ we get by \Cref{eq:complete_field_at_generic_point} 
     and \Cref{theorem:kummer_ramification}, that $\cP$ is tamely ramified at the generic point of the exceptional divisor of the blowup.
\end{proof}

So, we conclude 
\begin{cor}\label{cor:ramification_after_blowup_is_tame}
    In the general settings of a curve $C$, and $\cP \to U=C \setminus \mdls$ a $G$-torsor with ramification bounded by $\ndls = n P \subset \mdls$ at $P$, we have:
    \begin{enumerate}
        \item If $G=\Z/p\Z$ then $\cP^{[\deg \ndls]}$ is unramified at $\eta_\ndls$
        \item If $\cP$ is tamely ramified at $\ndls$ then $\cP^{[\deg \ndls]}$ is tamely ramified at $\eta_\ndls$
    \end{enumerate}
\end{cor}
\begin{proof}
    By the remark opening this subsection we may assume $k$ is algebraically closed, so that $P \in C(k)$ and
    $\deg \ndls = n$. The restriction $\cP_x$ of $\cP$ to the punctured formal disc at $P$ is tamely ramified, resp.\ a
    $\Z/p\Z$-torsor with ramification bounded by $n$, so by \Cref{lemma:realization_over_Gm} there is a torsor
    $\cP' \to \bG_m$ of the corresponding explicit form with $\cP'_x \cong \cP_x$. By \Cref{cor:reduction_to_P1}, the
    ramification of $\cP^{[\deg \ndls]}$ at $\eta_\ndls$ agrees with that of $\cP'^{[\deg \ndls]}$ at $\eta_{n \cdot 0}$,
    which was computed in \Cref{theorem:ramification_after_blowup_P1}.
\end{proof}

% =====================================================================
%  The case G = Z/p^r Z.
%  ---------------------------------------------------------------
%  Imported from the working note induction_completion/gap_closure
%  ("Closing the gaps"), adapted to the thesis conventions.  The
%  general Witt-vector background it relies on lives in the
%  Preliminaries (witt/WittVectorPrerequisites.tex); this file contains
%  only the parts specific to the theorem: the W_r-realization lemma
%  over G_m (Gap 1), the identification of the symmetric-power cover
%  (Gap 2), the degree count, and the proof itself.
% =====================================================================

Next, we generalize to the case $G = \Z/p^r\Z$.

\begin{theorem}\label{theorem:ramification_after_blowup_is_tame_for_p_power}
    Assume $G=\Z/p^r\Z$. Then $\cP^{[\deg (\ndls)]}$ is unramified at $\eta_{\ndls}$.
\end{theorem}

The proof occupies the remainder of this subsection. In contrast with what the phrase
``we generalize'' may suggest, the argument is \emph{not} an induction on $r$ with
\Cref{cor:ramification_after_blowup_is_tame}(1) as its base case: it is a direct
argument, treating all $r$ at once, in which each of the three steps of the
Artin--Schreier case of \Cref{theorem:ramification_after_blowup_P1} --- realization of
the local torsor by an explicit equation over $\bG_m$, identification of the
symmetric-power cover at $\eta_\ndls$, and the degree estimate --- is replaced by its
Witt-vector analogue, using the theory collected in
\Cref{subsection:witt_vector_prerequisites}. In particular the case $r=1$ is
\emph{reproved} (as the special case of Witt vectors of length one) rather than used.
For an explanation of why the seemingly natural induction on $r$ fails --- in short:
the hypothesis bounds the ramification of $\cP$ in \emph{upper} numbering, while the
break of the top degree-$p$ layer $\cP \to \cP'$ is the largest \emph{lower}-numbering
break, which can be of size about $p^{r-1}d$ --- see \Cref{section:a_wrong_path}.

\subsubsection*{Standing reduction}

As in the proof of \Cref{cor:ramification_after_blowup_is_tame}, we may assume $k$ is
algebraically closed. We may further assume that $\cP \to U$ is \emph{totally
ramified} at $P$ with inertia equal to all of $G$: writing
$R = \widehat{\Oo_{C,P}}$, $K = \Frac(R)$, and letting
$H = \rho(G^0) \subseteq G$ be the image of the inertia subgroup under the
homomorphism $\rho \colon \Gal(K^{sep}/K) \to G$ classifying $\cP|_{\Spec K}$,
decompose $\cP \to \cP_{G/H} \to U$ as in \Cref{prop:quotient_torsors}, where
$\cP \to \cP_{G/H}$ is an $H$-torsor and $\cP_{G/H} \to U$ a $G/H$-torsor. Then
$\cP_{G/H}$ is unramified at $P$, hence extends over $P$; taking any point
$Q \in \cP_{G/H}$ over $P$ yields $\widehat{\Oo_{\cP_{G/H}, Q}} \cong \widehat{\Oo_{C,P}}$,
and since $k$ is algebraically closed the torsor $\cP_{G/H}$ is locally trivial at
$Q$. This replaces $(G, \cP \to U)$ by $(H, \cP \to \cP_{G/H})$ with
$H \cong \Z/p^s\Z$ for some $s \leq r$; as the theorem is proved for all $r$
simultaneously, we may rename and assume $H = G$.

After this reduction, the restriction $\cQ := \cP_x$ of $\cP$ to the punctured formal
disc $\Spec k((x))$ at $P$ (with $x$ a uniformizer at $P$) is a \emph{connected}
$\Z/p^r\Z$-torsor: $\cQ = \Spec L$ for a cyclic, totally ramified extension $L/K$ of
degree $p^r$, with ramification bounded by $d := \deg \ndls$
(\Cref{definition:local_ramification_boundness}); equivalently, the largest
upper-numbering break $u$ of $L/K$ satisfies $u < d$ \emph{strictly}.

\subsubsection*{Realization over $\bG_m$ for Witt vectors}

The first step generalizes \Cref{lemma:realization_over_Gm}(2) from $\Z/p\Z$ to
$\Z/p^r\Z$: every local torsor class as above is realized by an explicit
Artin--Schreier--Witt equation over $\bG_m$ whose Witt datum is a vector of
\emph{polynomials in $x^{-1}$} with controlled pole orders. Recall from
\Cref{definition:reduced_witt_vector} the notion of a reduced vector; we add:

\begin{definition}\label{definition:reduced_polynomial_vector}
    A vector $\vec g = (g_0, \dots, g_{r-1}) \in W_r(k((x)))$ is a \textbf{reduced
    polynomial vector} if it is reduced (\Cref{definition:reduced_witt_vector}) and
    moreover every $g_j \in x^{-1}k[x^{-1}]$; then $m(g_j) = \deg_{x^{-1}} g_j$, and
    $m(g_j) = 0$ if and only if $g_j = 0$.
\end{definition}

\begin{lemma}[surjectivity of $\wp$ on integral Witt vectors]\label{lemma:wp_surjective_on_integral_witt}
    $\wp \colon W_r(k[[x]]) \to W_r(k[[x]])$ is surjective.
\end{lemma}

\begin{proof}
    Induction on $r$. For $r=1$ this is the first bullet in the proof of
    \Cref{lemma:realization_over_Gm}(2): given $h = \sum_{i \geq 0} b_i x^i$, solve
    $u^p - u = h$ coefficientwise, using that $k$ is algebraically closed for the
    constant term. For $r>1$, let $\vec h \in W_r(k[[x]])$. By induction choose
    $\vec u' \in W_{r-1}(k[[x]])$ with $\wp \vec u' = t(\vec h)$, and lift it to
    $\tilde u := (\vec u', 0) \in W_r(k[[x]])$. By
    \Cref{lemma:witt_bookkeeping}(i),
    $t\bigl(\wp \tilde u \ominus \vec h\bigr) = \wp \vec u' \ominus t(\vec h) = 0$, so
    $\wp \tilde u \ominus \vec h = V^{r-1}(g)$ for some $g \in k[[x]]$. Choose
    $w \in k[[x]]$ with $\wp w = -g$ (case $r=1$ again). Then by
    \Cref{lemma:witt_bookkeeping}(ii) and additivity of $V^{r-1}$,
    \[
    \wp\bigl(\tilde u \oplus V^{r-1}(w)\bigr)
    = \wp \tilde u \oplus V^{r-1}(\wp w)
    = \bigl(\vec h \oplus V^{r-1}(g)\bigr) \oplus V^{r-1}(-g) = \vec h. \qedhere
    \]
\end{proof}

\begin{prop}[reduced polynomial representatives]\label{prop:reduced_polynomial_representatives}
    Every class in $W_r(k((x)))/\wp W_r(k((x)))$ has a reduced polynomial
    representative: for every $\vec f \in W_r(k((x)))$ there exists
    $\vec u \in W_r(k((x)))$ such that $\vec g := \vec f \ominus \wp \vec u$ has all
    components in $x^{-1}k[x^{-1}]$ with pole orders either $0$ or prime to $p$.
\end{prop}

\begin{proof}
    Induction on $r$.

    \emph{Base case $r=1$.} This is the argument of
    \Cref{lemma:realization_over_Gm}(2): first, writing $f = f^- + f^+$ with
    $f^- \in x^{-1}k[x^{-1}]$ and $f^+ \in k[[x]]$, surjectivity of $\wp$ on $k[[x]]$
    removes $f^+$; then, as long as the leading term is $cx^{-m}$ with $p \mid m$,
    $m>0$, subtracting $\wp(c^{1/p}x^{-m/p}) = cx^{-m} - c^{1/p}x^{-m/p}$ strictly
    decreases the pole order while staying inside $x^{-1}k[x^{-1}]$; this terminates
    at $0$ or at pole order prime to $p$.

    \emph{Inductive step.} Let $\vec f \in W_r(k((x)))$. By induction there is
    $\vec u' \in W_{r-1}(k((x)))$ such that
    $\vec g' := t(\vec f) \ominus \wp \vec u'$ is a reduced polynomial vector in
    $W_{r-1}$. Lift $\vec u'$ to $\tilde u := (\vec u', 0) \in W_r(k((x)))$ and set
    $\vec h := \vec f \ominus \wp \tilde u$. By \Cref{lemma:witt_bookkeeping}(i),
    $t(\vec h) = \vec g'$, i.e.\
    \[
    \vec h = (g'_0, \dots, g'_{r-2},\, h_{r-1}), \qquad h_{r-1} \in k((x)),
    \]
    with the first $r-1$ components already reduced polynomials. It remains to repair
    the last component \emph{without touching the others}; by
    \Cref{lemma:witt_bookkeeping}(iii), adjustments of the form
    $\ominus\, \wp\bigl(V^{r-1}(z)\bigr) = \ominus\, V^{r-1}(\wp z)$ change only the
    last component, by $-\wp(z)$. So:
    \begin{itemize}
        \item Write $h_{r-1} = h^- + h^+$ with $h^- \in x^{-1}k[x^{-1}]$,
        $h^+ \in k[[x]]$, choose $w \in k[[x]]$ with $\wp w = h^+$
        (\Cref{lemma:wp_surjective_on_integral_witt} with $r=1$), and replace
        $\vec h$ by $\vec h \ominus \wp(V^{r-1}w)$: the last component becomes $h^-$.
        \item While the last component has leading term $cx^{-m}$ with $p \mid m$,
        $m>0$, replace $\vec h$ by
        $\vec h \ominus \wp\bigl(V^{r-1}(c^{1/p}x^{-m/p})\bigr)$: the last component
        changes by $-\wp(c^{1/p}x^{-m/p}) = -cx^{-m} + c^{1/p}x^{-m/p}$, so its pole
        order strictly decreases and it stays in $x^{-1}k[x^{-1}]$. This terminates at
        $0$ or at pole order prime to $p$.
    \end{itemize}
    The result is a reduced polynomial vector $\wp$-equivalent to $\vec f$.
\end{proof}

\begin{rem*}
    The individual components of a Witt vector cannot be truncated to their principal
    parts one at a time --- Witt addition is not componentwise
    (\Cref{eq:witt_addition_shape}), so a componentwise operation changes the
    Artin--Schreier--Witt class. The point of
    \Cref{prop:reduced_polynomial_representatives} is that all the adjustments are
    performed by subtracting genuine elements of $\wp W_r(k((x)))$, inside the Witt
    ring.
\end{rem*}

\begin{lemma}[Realization over $\bG_m$ for $W_r$]\label{lemma:realization_over_Gm_Wr}
    Let $k$ be algebraically closed of characteristic $p>0$ and let $\cQ$ be a
    \emph{connected} $\Z/p^r\Z$-torsor over $\Spec(k((x)))$ with ramification bounded
    by $d$ (\Cref{definition:local_ramification_boundness}). Then there is a reduced
    polynomial vector $\vec f = (f_0, \dots, f_{r-1})$, $f_j \in x^{-1}k[x^{-1}]$,
    such that, with $m_j := \deg_{x^{-1}} f_j$:
    \begin{enumerate}
        \item \label{item:realization_Wr_polebound} $p^{\,r-1-j}\, m_j < d$ for all
        $0 \leq j \leq r-1$; equivalently $m_j < d\, p^{\,j-(r-1)}$;
        \item \label{item:realization_Wr_torsor}
        $\cP' := \Spec\Bigl(k[x,x^{-1}][\vec X]\big/\bigl(\wp \vec X \ominus \vec f\bigr)\Bigr)$
        is a $\Z/p^r\Z$-torsor over $\bG_m$ with $\cP'_x \cong \cQ$;
        \item \label{item:realization_Wr_infty} $\cP'$ extends to a torsor over
        $\bP^1_k \setminus \{0\}$ --- in particular it is unramified at $\infty$ ---
        and has ramification bounded by $d$ at $0$.
    \end{enumerate}
\end{lemma}

\begin{proof}
    $\cQ = \Spec L$ with $L/K$ cyclic totally ramified of degree $p^r$ and largest
    upper break $u < d$ (\Cref{theorem:asw_theory}). Let
    $[\vec f_0] \in W_r(K)/\wp$ be its Artin--Schreier--Witt class and let $\vec f$ be
    a reduced polynomial representative
    (\Cref{prop:reduced_polynomial_representatives}).

    First, $f_0 \neq 0$: otherwise
    \Cref{cor:asw_class_order_and_first_component} would bound the order of the class
    by $p^{r-1}$, contradicting $[L:K] = p^r$. Hence $m_0 \geq 1$ and $p \nmid m_0$,
    and the Schmid--Witt break formula
    (\Cref{theorem:schmid_witt_break_formula}) applies:
    $\max_j p^{r-1-j} m_j = u < d$, which is \ref{item:realization_Wr_polebound}.

    For \ref{item:realization_Wr_torsor}: by \Cref{lemma:asw_etale_torsor} (with
    $A = k[x,x^{-1}] \ni f_j$), $\cP'$ is a $\Z/p^r\Z$-torsor over $\bG_m$; its
    restriction to $\Spec k((x))$ is the Artin--Schreier--Witt torsor of the class of
    $\vec f$, i.e.\ $\cQ$ (\Cref{theorem:asw_theory}).

    For \ref{item:realization_Wr_infty}: in the coordinate $y = x^{-1}$ at $\infty$ we
    have $f_j \in y\,k[y]$, so $\vec f \in W_r(k[y])$ and
    \Cref{lemma:asw_etale_torsor} (with $A = k[y]$) extends $\cP'$ to a torsor over
    $\A^1_k = \Spec k[y] = \bP^1_k \setminus \{0\}$; a torsor over the full local ring
    at $\infty$ is unramified there. The ramification bound at $0$ holds because
    $\cP'_x \cong \cQ$.
\end{proof}

\begin{rem*}\label{rem:realization_Wr_sanity}
    Connectedness of $\cQ$ is guaranteed in the application by the standing reduction
    (inertia equal to all of $G$). Note also the following sanity check: for every
    $j$ with $p^{\,r-1-j} \geq d$, the bound \ref{item:realization_Wr_polebound}
    forces $m_j = 0$, i.e.\ $f_j = 0$ --- all sufficiently low Witt components of a
    reduced representative vanish outright. This makes visible how strong the
    hypothesis $u < d$ is at the lower levels.
\end{rem*}

\subsubsection*{The symmetric-power cover at $\eta_\ndls$}

We now generalize the invariant-ring computation in the Artin--Schreier case of
\Cref{theorem:ramification_after_blowup_P1} to Witt coordinates. Fix
$\vec f \in W_r(R)$, $R = k[x, x^{-1}]$, and let
\[
S \;=\; R[\vec X]\big/\bigl(\wp \vec X \ominus \vec f\bigr)
\]
be the coordinate ring of the torsor $\cP' \to \bG_m$ of
\Cref{lemma:realization_over_Gm_Wr}. Write
$R^{\otimes_k d} = k[x_1^{\pm 1}, \dots, x_d^{\pm 1}]$ and
\[
S^{\otimes_k d}
= R^{\otimes_k d}[\vec X_1, \dots, \vec X_d]\big/
\bigl(\wp \vec X_1 \ominus \vec f(x_1), \dots, \wp \vec X_d \ominus \vec f(x_d)\bigr),
\]
where $\vec X_i = (X_{i,0}, \dots, X_{i,r-1})$ and $\vec f(x_i)$ is the image of
$\vec f$ under $x \mapsto x_i$. Exactly as in the proof of
\Cref{theorem:ramification_after_blowup_P1}, the $d$-fold product torsor
$p_1^{-1}\cP' \otimes^G \cdots \otimes^G p_d^{-1}\cP'$ on $U^d$ is the quotient of
$\Spec S^{\otimes_k d}$ by $\KG \cong G^{d-1}$, embedded in $G^d$ as the kernel of the
sum map via $(g_1, \dots, g_{d-1}) \mapsto (g_1, g_2 - g_1, \dots, -g_{d-1})$, and
acting (in Witt coordinates, writing $\underline g$ for the image of $g \in G$ under
$\Z/p^r\Z \cong W_r(\F_p)$, \Cref{example:witt_of_Fp}) by
\[
(g_1, \dots, g_{d-1}) \colon \quad
\vec X_1 \mapsto \vec X_1 \oplus \underline g_1, \qquad
\vec X_i \mapsto \vec X_i \oplus \underline g_i \ominus \underline g_{i-1}
\ (1 < i < d), \qquad
\vec X_d \mapsto \vec X_d \ominus \underline g_{d-1}.
\]

\begin{prop}[identification of the product torsor]\label{prop:witt_cover_identification}
    Set
    \[
    \vec Y := \bigoplus_{i=1}^{d} \vec X_i \in W_r\bigl(S^{\otimes_k d}\bigr),
    \qquad
    \vec \alpha := \bigoplus_{i=1}^{d} \vec f(x_i) \in W_r\bigl(R^{\otimes_k d}\bigr),
    \]
    where $\bigoplus$ denotes the $d$-fold Witt vector sum $\oplus$ (not a direct
    sum); in particular $\vec Y$ is a single Witt vector of length $r$. Then:
    \begin{enumerate}
        \item \label{item:witt_cover_invariance} Every component $Y_n$ of $\vec Y$ is
        $G^{d-1}$-invariant, and
        \[
        \bigl(S^{\otimes_k d}\bigr)^{G^{d-1}}
        \;\cong\;
        R^{\otimes_k d}[\vec Y]\big/\bigl(\wp \vec Y \ominus \vec \alpha\bigr)
        \]
        as $G$-torsors over $R^{\otimes_k d}$ (the residual $G = G^d/G^{d-1}$ acting
        by $\vec Y \mapsto \vec Y \oplus \underline g$).
        \item \label{item:witt_cover_symmetry} Every component $\alpha_n$ of
        $\vec \alpha$ is a symmetric Laurent polynomial, i.e.\ lies in
        $\bigl(R^{\otimes_k d}\bigr)^{S_d} = k[e_1, \dots, e_d, e_d^{-1}]$, and
        \[
        \Bigl(\bigl(S^{\otimes_k d}\bigr)^{G^{d-1}}\Bigr)^{S_d}
        \;\cong\;
        k[e_1, \dots, e_d, e_d^{-1}][\vec Y]\big/\bigl(\wp \vec Y \ominus \vec \alpha\bigr).
        \]
    \end{enumerate}
    Consequently, restricting to the complete local field
    $\widehat K_\eta = \kappa(\eta)((e_d))$,
    $\kappa(\eta) = k(e_1/e_d, \dots, e_{d-1}/e_d)$, of
    \Cref{eq:complete_field_at_generic_point},
    \[
    \cP'^{[d]}\big|_{\Spec \widehat K_\eta}
    \;\cong\;
    \Spec\; \widehat K_\eta[\vec Z]\big/\bigl(\wp \vec Z \ominus \vec \alpha\bigr).
    \]
\end{prop}

\begin{proof}
    \ref{item:witt_cover_invariance} An element
    $\sigma = (g_1, \dots, g_{d-1}) \in G^{d-1}$ acts as a ring automorphism of
    $S^{\otimes_k d}$: this is its \emph{original} action, defined on the generators
    $X_{i,n}$ by the twisted formulas above (which mix components through the carry
    polynomials of \Cref{eq:witt_addition_shape}). By functoriality of $W_r$
    (\Cref{subsection:witt_vector_prerequisites}), the induced map $W_r(\sigma)$ on
    $W_r(S^{\otimes_k d})$ is componentwise and commutes with $\oplus, \ominus$; on
    the generators it reproduces the twisted action,
    $W_r(\sigma)(\vec X_i) = \vec X_i \oplus \underline g_i \ominus \underline g_{i-1}$
    (with $\underline g_0 = \underline g_d = \vec 0$). Hence, using commutativity and
    associativity of $\oplus$, the Witt sum transforms by the telescoping constant
    \[
    W_r(\sigma)(\vec Y)
    \;=\; \bigoplus_{i=1}^{d} W_r(\sigma)(\vec X_i)
    \;=\; \vec Y \oplus \underline g_1 \oplus (\underline g_2 \ominus \underline g_1)
    \oplus \cdots \oplus (\ominus\, \underline g_{d-1})
    \;=\; \vec Y.
    \]
    Since $W_r(\sigma)$ is componentwise, this equality of $r$-tuples says exactly
    $\sigma(Y_n) = Y_n$ for every $n$, i.e.\ each $Y_n \in S^{\otimes_k d}$ is
    invariant. Since $\wp = F - \mathrm{id}$ is additive
    (\Cref{definition:asw_operator} ff.),
    \[
    \wp \vec Y = \bigoplus_{i=1}^{d} \wp \vec X_i
    = \bigoplus_{i=1}^{d} \vec f(x_i) = \vec \alpha .
    \]
    Hence there is a well-defined $R^{\otimes_k d}$-algebra homomorphism
    \[
    \phi \colon\; T := R^{\otimes_k d}[\vec Z]\big/\bigl(\wp \vec Z \ominus \vec \alpha\bigr)
    \;\longrightarrow\; \bigl(S^{\otimes_k d}\bigr)^{G^{d-1}},
    \qquad \vec Z \mapsto \vec Y .
    \]
    By \Cref{lemma:asw_etale_torsor}, $\Spec T$ is a $G$-torsor over $U^d$. The
    target is the product torsor (the contracted product of the $p_i^{-1}\cP'$ along
    the sum map $G^d \to G$), a $G$-torsor over $U^d$ by
    \Cref{prop:symmetric_power_torsor} and the surrounding discussion; its residual
    $G$-action is induced by $(g, 0, \dots, 0) \in G^d$, i.e.\
    $\vec X_1 \mapsto \vec X_1 \oplus \underline g$, under which
    $\vec Y \mapsto \vec Y \oplus \underline g$. So $\Spec \phi$ is a $G$-equivariant
    morphism of $G$-torsors over $U^d$, hence an isomorphism by
    \Cref{prop:torsor_morphism_iso}.

    \ref{item:witt_cover_symmetry} $S_d$ acts on $S^{\otimes_k d}$ by permuting the
    tensor factors, i.e.\ $x_i \mapsto x_{\sigma(i)}$,
    $\vec X_i \mapsto \vec X_{\sigma(i)}$. As in
    \ref{item:witt_cover_invariance}, each $\sigma \in S_d$ is a ring automorphism,
    so $W_r(\sigma)$ is componentwise and commutes with $\oplus$. Since $\oplus$ is
    commutative and associative, permuting the summands of a Witt sum leaves it
    unchanged; thus $\vec Y$ is $S_d$-invariant (componentwise, as before) and
    $W_r(\sigma)(\vec \alpha) = \bigoplus_i \vec f(x_{\sigma(i)}) = \vec \alpha$, so
    each $\alpha_n \in R^{\otimes_k d}$ is a symmetric Laurent polynomial, hence lies
    in $k[e_1, \dots, e_d, e_d^{-1}]$ (the $S_d$-invariants of $R^{\otimes_k d}$, as
    in the proof of \Cref{theorem:ramification_after_blowup_P1}). By
    \Cref{lemma:asw_etale_torsor}, $T$ is free over $R^{\otimes_k d}$ with basis the
    monomials $\prod_n Y_n^{a_n}$, $0 \leq a_n < p$, each of which is
    $S_d$-invariant; writing an element in this basis, it is $S_d$-invariant if and
    only if all its coefficients are. Therefore
    \[
    T^{S_d} = \bigl(R^{\otimes_k d}\bigr)^{S_d}[\vec Y]\big/\bigl(\wp \vec Y \ominus \vec \alpha\bigr)
    = k[e_1, \dots, e_d, e_d^{-1}][\vec Y]\big/\bigl(\wp \vec Y \ominus \vec \alpha\bigr).
    \]

    The final statement follows exactly as in the $r=1$ case: by definition
    $\cP'^{[d]}$ is the quotient by $\Xi = \KG \rtimes S_d$ computed above
    (\Cref{cor:sym-power-quotient}), and its restriction to $\Spec \widehat K_\eta$
    along $k[e_1, \dots, e_d, e_d^{-1}] \to \widehat K_\eta$ is the
    Artin--Schreier--Witt cover of the image of $\vec \alpha$ in
    $W_r(\widehat K_\eta)$.
\end{proof}

\begin{rem*}
    The proof via \Cref{prop:torsor_morphism_iso} replaces the rank count of the
    $r=1$ argument ($p^{rd}/p^{r(d-1)} = p^r$ together with irreducibility of the
    invariant's minimal equation); the rank count alone would still require showing
    that $\vec Y$ generates the invariant ring, which is exactly what equivariance
    gives for free. Note also that a pure restriction argument on $H^1$'s (``the
    class of the product torsor is the sum of the pulled-back characters'')
    identifies the cover over $k(x_1, \dots, x_d)$ but does \emph{not} by itself
    descend it to the symmetric base $k(e_1, \dots, e_d)$; the invariant-theoretic
    computation above is what effects the descent.
\end{rem*}

\subsubsection*{The degree count}

Throughout this paragraph, $\vec f$ is the reduced polynomial vector of
\Cref{lemma:realization_over_Gm_Wr}, so $m_j = \deg_{x^{-1}} f_j$ satisfies
$m_j < d\, p^{\,j-(r-1)}$, and $y_i := 1/x_i$, so that $f_j(x_i) \in y_i\, k[y_i]$ is
a polynomial in $y_i$ of degree $m_j$.

\begin{lemma}[symmetric regularity]\label{lemma:symmetric_regularity}
    Let $\Phi \in k[y_1, \dots, y_d]^{S_d}$ be symmetric with every monomial of total
    degree $< d$. Then $\Phi \in \kappa(\eta) = k(e_1/e_d, \dots, e_{d-1}/e_d)$; in
    particular $\Phi$ is regular at $\eta_\ndls$.
\end{lemma}

\begin{proof}
    $\Phi$ is a $k$-combination of monomial symmetric functions $m_\lambda(y)$ with
    $|\lambda| < d$. Each $m_\lambda(y)$ is an isobaric polynomial in the elementary
    symmetric functions $e_\mu(y) = \prod_i e_{\mu_i}(y)$ with $\mu \vdash |\lambda|$;
    every part $\mu_i \leq |\lambda| < d$, so only $e_1(y), \dots, e_{d-1}(y)$ occur
    --- never $e_d(y)$. Now $e_k(y) = e_{d-k}(x)/e_d(x)$, which for
    $1 \leq k \leq d-1$ lies in $\kappa(\eta)$. Hence
    $\Phi \in \kappa(\eta) \subset \widehat\Oo_\eta = \kappa(\eta)[[e_d]]$.
\end{proof}

\begin{theorem}[regularity of all Witt components]\label{theorem:witt_components_regularity}
    With $\vec f$ as in \Cref{lemma:realization_over_Gm_Wr}, every component
    $\alpha_n$ $(0 \leq n \leq r-1)$ of
    $\vec \alpha = \bigoplus_{i=1}^{d} \vec f(x_i)$ lies in $\kappa(\eta)$; in
    particular $\vec \alpha \in W_r\bigl(\widehat\Oo_\eta\bigr)$.
\end{theorem}

\begin{proof}
    Fix $n \leq r-1$ and a monomial $\prod_{i,j} \bigl(f_j(x_i)\bigr)^{e_{ij}}$ of
    $\alpha_n$, viewed as a universal isobaric polynomial in the $f_j(x_i)$ of
    weight $p^n$ (\Cref{lemma:witt_addition_isobaric}, with $f_j(x_i)$ of weight
    $p^j$); set $E_j := \sum_i e_{ij}$, so $\sum_j E_j p^{\,j} = p^{\,n}$. Expanding
    each $f_j(x_i)$ as a polynomial in $y_i$ of degree $m_j$, every resulting
    monomial in the $y_i$ has total degree
    \[
    \leq\; \sum_{i,j} e_{ij}\, m_j \;=\; \sum_j E_j\, m_j
    \;<\; \sum_j E_j\, d\, p^{\,j-(r-1)}
    \;=\; d\, p^{-(r-1)} \sum_j E_j p^{\,j}
    \;=\; d\, p^{\,n-(r-1)} \;\leq\; d,
    \]
    strictly: some $E_j > 0$ (as $\sum_j E_j p^j = p^n > 0$), and the pole bound
    $m_j < d\, p^{\,j-(r-1)}$ of
    \Cref{lemma:realization_over_Gm_Wr}\ref{item:realization_Wr_polebound} is strict
    even when $m_j = 0$. Thus every monomial of $\alpha_n$ --- a symmetric polynomial
    in $y_1, \dots, y_d$ by
    \Cref{prop:witt_cover_identification}\ref{item:witt_cover_symmetry} --- has total
    $y$-degree $< d$, and \Cref{lemma:symmetric_regularity} gives
    $\alpha_n \in \kappa(\eta)$.
\end{proof}

\begin{rem*}
    The bound is exact: the weight $p^{\,j-(r-1)}$ is paired against the isobaricity
    constraint $\sum_j E_j p^{\,j} = p^{\,n}$, so the estimate
    $d\, p^{\,n-(r-1)} \leq d$ is independent of how the Witt carries distribute the
    exponents $E_j$; the hypothesis ``break $< d$ at $P$'' is used at \emph{every}
    Witt level at once. The sanity check of
    \Cref{lemma:realization_over_Gm_Wr} records the extreme case: components with
    $p^{\,r-1-j} \geq d$ vanish outright in the reduced representative.
\end{rem*}

\subsubsection*{Proof of the theorem}

\begin{proof}[Proof of \Cref{theorem:ramification_after_blowup_is_tame_for_p_power}]
    By the standing reduction above we may assume $k$ algebraically closed and
    $\cP \to U$ totally ramified at $P$ with inertia all of $G$, so that
    $\cQ := \cP_x$ is a connected $\Z/p^r\Z$-torsor over $\Spec k((x))$ with
    ramification bounded by $d = \deg \ndls$.

    \Cref{lemma:realization_over_Gm_Wr} produces a torsor $\cP' \to \bG_m$, defined
    by a reduced polynomial vector $\vec f$ with $p^{r-1-j} m_j < d$, with
    $\cP'_x \cong \cQ$. By \Cref{cor:reduction_to_P1} (valid for arbitrary finite
    abelian $G$), the ramification of $\cP^{[d]}$ at $\eta_\ndls$ agrees with that of
    $\cP'^{[d]}$ at $\eta_{d \cdot 0}$; so we may assume $C = \bP^1_k$,
    $\mdls = d \cdot 0 + \infty$, $\ndls = d \cdot 0$, $\cP = \cP'$.

    By \Cref{prop:witt_cover_identification},
    \[
    \cP'^{[d]}\big|_{\Spec \widehat K_\eta}
    \;\cong\;
    \Spec\; \widehat K_\eta[\vec Z]\big/\bigl(\wp \vec Z \ominus \vec \alpha\bigr),
    \qquad
    \vec \alpha = \bigoplus_{i=1}^{d} \vec f(x_i),
    \]
    with $\widehat K_\eta = \kappa(\eta)((e_d))$ the complete local field at
    $\eta_\ndls$ (\Cref{eq:complete_field_at_generic_point}) and
    $\widehat\Oo_\eta = \kappa(\eta)[[e_d]]$ its valuation ring.

    By \Cref{theorem:witt_components_regularity},
    $\vec \alpha \in W_r(\kappa(\eta)) \subset W_r(\widehat\Oo_\eta)$. Hence, by
    \Cref{lemma:asw_etale_torsor} applied over $A = \widehat\Oo_\eta$, the
    Artin--Schreier--Witt cover $\wp \vec Z = \vec \alpha$ is finite \'etale over
    $\widehat\Oo_\eta$, and $\cP'^{[d]}\big|_{\Spec \widehat K_\eta}$ is its generic
    fibre: every field factor of
    $\widehat K_\eta[\vec Z]/(\wp \vec Z \ominus \vec \alpha)$ is the fraction field
    of a finite \'etale $\widehat\Oo_\eta$-algebra, i.e.\ an unramified extension of
    $\widehat K_\eta$ with respect to the discrete valuation ring
    $\widehat\Oo_\eta$. By \Cref{definition:tamely_ramified_in_codim_1},
    $\cP^{[d]}$ is unramified at $\eta_\ndls$.
\end{proof}

\begin{rem*}
    The argument proves all $r$ simultaneously; it reproves the case $r=1$
    (\Cref{cor:ramification_after_blowup_is_tame}(1)) as the special case of a
    vector of length one and never invokes it.
\end{rem*}


Finally, \textbf{Proof of Theorem \cref{theorem:torsors_after_blowup_is_tame}}:
Idea: decompose $G$ be the structure theorem, and finish by the above. 



\subsection{Extending To Product of Blowups}\label{subsection:ram_prod_analysis}
We conclude this sesction by proving sume auxilary results that would be useful in \Cref{section:gcft}.


%%% %%% %%% %%% %%% %%% %%% %%% %%% %%% %%% %%% %%% %%% %%% %%%
%       Behavior of Ramification under Product of Blowups
%%% %%% %%% %%% %%% %%% %%% %%% %%% %%% %%% %%% %%% %%% %%% %%%
%\subsubsection{Behavior of Ramification under Product of Blowups}

Let $X$ and $Y$ be smooth schemes over a field $k$, and let $x \in X$ and $y \in Y$ be closed points. 
We denote the blowups of these schemes at the given points by 
$\pi_X: \Bl_x(X) \to X$ and $\pi_Y: \Bl_y(Y) \to Y$. Furthermore, let $\pi_{X \times Y}: \Bl_{(x,y)}(X \times_k Y) \to X \times_k Y$ 
be the blowup of the product scheme at the point $(x,y)$.
We denote by $E_X, E_Y$, and $E_{X \times Y}$ the respective exceptional divisors, and let $\eta_X, \eta_Y$, and $\eta_{X \times Y}$ be their generic points.

In this section, we establish the following result concerning the stability of ramification bounds under the external product of torsors.

\begin{prop}\label{prop:ramification-external-product}
Let $G$ be a finite abelian group. Suppose $\cG_X$ and $\cG_Y$ are $G$-torsors defined on open subsets $U_X \subset \Bl_x(X)$ and $U_Y \subset \Bl_y(Y)$ that are disjoint 
from the exceptional divisors. 
If the ramification of $\cG_X$ at $\eta_X$ and $\cG_Y$ at $\eta_Y$ is bounded by $r$, then the external product torsors 
\[
\cG_{X \times Y} := \text{pr}_1^{-1} \cG_X \otimes \text{pr}_2^{-1} \cG_Y
\]
has ramification bounded by $r$ at the generic point $\eta_{X\times Y}$ of the exceptional divisor in the product blowup.
\end{prop}

The proposition is purely local in nature, it suffices to consider the case where $X$ and $Y$ are affine. 
More precisely, by the smoothness of $X$ and $Y$, we may restrict our attention to open neighborhoods of $x$ and $y$ that are etale over affine spaces. 
The rest of this section treats that case.

\subsubsection*{The Affine Case}
Let $X = \mathbb{A}_k^n$ be the affine $n$-space over a field $k$, and let $0 \in X$ be the origin. 
Let $\tilde{X} = \text{Bl}_0(X)$ be the blowup of $X$ at the origin. 
Recall that $\tilde{X} \subset X \times_k \mathbb{P}^{n-1}_k$ is defined by the equations $x_i u_j = x_j u_i$, where $[u_1 : \dots : u_n]$ are the homogeneous coordinates of 
$\mathbb{P}^{n-1}_k$.
The exceptional divisor $E \subset \tilde{X}$ is the fiber over the origin, $E = \{ (0, [u_1 : \dots : u_n]) \}$, which is of codimension 1 in $\tilde{X}$. 
Let $\eta \in E$ be the generic point of $E$, and let $R = \mathcal{O}_{\tilde{X}, \eta}$ be the associated local ring. 
This ring $R$ is a discrete valuation ring (DVR) with fraction field $K = K(X) = k(x_1, \dots, x_n)$.

On the affine chart $U_1$ where $u_1 \neq 0$, we have $x_i = \frac{u_i}{u_1} x_1$. 
The coordinate ring is:$$\mathcal{O}_{\tilde{X}}(U_1) = k\left[x_1, \frac{u_2}{u_1}, \dots, \frac{u_n}{u_1}\right]$$
In this chart, the generic point $\eta$ corresponds to the prime ideal $\mathfrak{p}_1 = (x_1)$. Thus, the local ring is $R = k[x_1, \frac{u_2}{u_1}, \dots, \frac{u_n}{u_1}]_{(x_1)}$. 
The residue field is $\kappa(\eta) = k(\frac{u_2}{u_1}, \dots, \frac{u_n}{u_1})$, and the completion of $R$ with respect to its maximal ideal is:
$$\hat{R} = \kappa(\eta)[[x_1]]$$
In this local ring, $x_1$ is a uniformizer. Note that any $x_i$ (for $i > 1$) can also serve as a uniformizer, 
as $x_i = (\frac{u_i}{u_1}) x_1$ and $\frac{u_i}{u_1}$ is a unit in $R$.

For any monomial $M = x_1^{a_1} \dots x_n^{a_n} \in K$, we can write:
$$M = x_1^{\sum a_i} \left( \frac{u_2}{u_1} \right)^{a_2} \dots \left( \frac{u_n}{u_1} \right)^{a_n}$$
Since the term in parentheses is a unit in $R$, the valuation $\nu_E$ associated with $E$ satisfies:
$$\nu_E(M) = \sum a_i = \deg M$$
Consequently, for any polynomial $f = f_d + f_{d+1} + \dots + f_l$, where $f_i$ is the homogeneous part of degree $i$, we have $\nu_E(f) = d$ (the order of vanishing at the origin).

\subsubsection*{The Product Case}
Now, let $X = \mathbb{A}^n$ and $Y = \mathbb{A}^m$ with origins $x=0$ and $y=0$. 
As before, $\text{Bl}_0(X) \subset X \times \mathbb{P}^{n-1}$ and $\text{Bl}_0(Y) \subset Y \times \mathbb{P}^{m-1}$ have exceptional divisors $E_X$ and $E_Y$ respectively.
Consider the product $X \times_k Y \cong \mathbb{A}_k^{n+m}$. 
The blowup of the product at the origin $(0,0)$, denoted $\text{Bl}_{(0,0)}(X \times_k Y)$, is a subscheme of $(X \times Y) \times \mathbb{P}^{n+m-1}$ defined by:

$$\begin{cases} 
x_i w_j = x_j w_i & 1 \le i, j \le n \\
y_k w_{n+l} = y_l w_{n+k} & 1 \le k, l \le m \\
x_i w_{n+j} = y_j w_i & 1 \le i \le n, \,\, 1 \le j \le m 
\end{cases}$$
where $[w_1 : \dots : w_{n+m}]$ are the homogeneous coordinates of $\mathbb{P}^{n+m-1}$. The exceptional divisor $E_{X \times Y}$ is isomorphic to $\mathbb{P}^{n+m-1}$.

\subsubsection*{Comparison of Blowups}

Both $\text{Bl}_0(X) \times_k \text{Bl}_0(Y)$ and $\text{Bl}_{(0,0)}(X \times_k Y)$ are birational to $X \times_k Y$. 
They share a common dense open set $\tilde{U}$ defined by the condition that neither the $X$-coordinates nor the $Y$-coordinates vanish simultaneously in the projective space:
$$\tilde{U} = \{ ((x,y), [w_1: \dots : w_{n+m}]) \mid (w_1, \dots, w_n) \neq 0 \text{ and } (w_{n+1}, \dots, w_{n+m}) \neq 0 \}$$This yields a diagram of open immersions:
$$ 
   \begin{tikzcd}
        & \tilde{U} \arrow[dl, "f_1"'] \arrow[dr, "f_2", hook] & \\
        \text{Bl}_0(X) \times_k \text{Bl}_0(Y) & & \text{Bl}_{(0,0)}(X \times_k Y)
    \end{tikzcd}
$$
    
    where $f_1$ maps the coordinates to the respective projectivizations $[w_1: \dots : w_n]$ and $[w_{n+1}: \dots : w_{n+m}]$.
    Also note that the generic point $\eta_{X \times Y}$  of $E_{X \times_k Y}$ is inside $\tilde{U}$

\subsubsection*{Extensions of DVRs}
Let $S$ be the local ring of the generic point $\eta_{X\times Y}$ of $E_{X \times Y}$ in $\tilde{U}$. 
On the chart where $w_1 \neq 0$ and $w_{n+1} \neq 0$, we have 
$x_1 = (\frac{w_1}{w_{n+1}}) y_1$. 
Since $\frac{w_1}{w_{n+1}}$ is a unit in this chart, $x_1$ and $y_1$ are equivalent as uniformizers.
We have:
$$
S=k\left[x_1, \frac{w_2}{w_1} \cdots, \frac{w_n}{w_1}, \frac{w_{n+1}}{w_1} \cdots \frac{w_{n+m}}{w_1}\right]_{(x_1)} = 
k\left[\frac{w_1}{w_{n+1}}, \frac{w_2}{w_{n+1}} \cdots \frac{w_n}{w_{n+1}}, y_1 \cdots \frac{w_{n+m}}{w_{n+1}}\right]_{(y_1)}
$$

$$k(\eta_{X \times Y}) = k\left(\frac{w_2}{w_1} \cdots, \frac{w_n}{w_1}, \frac{w_{n+1}}{w_1} \cdots \frac{w_{n+m}}{w_1}\right)$$

Let $R_X$ be the local ring of the exceptional divisor $E_X$ in $\text{Bl}_0(X)$. 
The pullback of $E_X \times_k \text{Bl}_0(Y)$ along $f_1$ induces an extension of DVRs $R_X \hookrightarrow S$. Which is:
\begin{enumerate}
    \item Weakly Unramified: $x_1$ is a uniformizer in both $R_X$ and $S$, so the ramification index is $e=1$.
    \item Residually Transcendental: The residue field extension $\kappa(\eta_X) \subset \kappa(\eta)$ is:
    $$k\left(\frac{w_2}{w_1}, \dots, \frac{w_n}{w_1}\right) \subset k\left(\frac{w_2}{w_1}, \dots, \frac{w_n}{w_1}, \frac{w_{n+1}}{w_1}, \dots, \frac{w_{n+m}}{w_1}\right)$$
    Hence separable.
\end{enumerate}
Since this extension is generated by transcendental elements, it is separable and formally smooth at the maximal ideal (\stackstag{09E7}).

\subsubsection*{Ramification of $G$-Torsors}
Let $G$ be a finite abelian group. 
Let $\mathcal{Q}$ be a $G$-torsor on an open $U \subset \text{Bl}_0(X)$ disjoint from $E_X$. 
%Suppose $\mathcal{Q}$ has ramification bounded by $r$ at the generic point $\eta_X$ of $E_X$.
Let $V \subset \text{Bl}_0(Y)$ be an open subscheme, and let $\pi_X: \text{Bl}_0(X) \times_k V \to \text{Bl}_0(X)$ 
be the projection onto the first factor. 
By restricting this projection to $U \times_k V$, we obtain the pullback $G$-torsor:
$$\pi_X^{-1}(\mathcal{Q}) \cong \mathcal{Q} \times_k V$$
which is defined on the open subset $U \times_k V \subset \text{Bl}_0(X) \times_k \text{Bl}_0(Y)$.
The extension of local rings $R_X \to S$ is weakly unramified (the ramification index $e=1$) and residually transcendental with seprable residue field exntesion. 
Under these conditions the ramification filtration is preserved. 
Therefore, the pullback $\pi_X^{-1}(\mathcal{Q})$ has ramification bounded by $r$ at the generic point $\eta_{X \times Y}$ of the exceptional divisor $E_{X \times Y}$ if and only if 
the original torsor $\mathcal{Q}$ has ramification bounded by $r$ at the generic point $\eta_X$ of $E_X$

And we finish by \Cref{lemma:bounding_ramification_of_contracted_product}.



